% Selector de bloques para la edición secuencial. La interfaz de llamada es
% \def\HAFOCaputoBlockSelect{1} junto con exactamente una de las banderas
% \HAFOCaputoBlockMemory, \HAFOCaputoBlockBasins,
% \HAFOCaputoBlockReturn o \HAFOCaputoBlockProtocol. Sin selector se conserva
% la sección completa y el mismo nivel de encabezados que en la versión original.
\def\HAFOCaputoShowMemory{1}
\def\HAFOCaputoShowBasins{1}
\def\HAFOCaputoShowReturn{1}
\def\HAFOCaputoShowProtocol{1}
\ifdefined\HAFOCaputoBlockSelect
  \def\HAFOCaputoShowMemory{0}
  \def\HAFOCaputoShowBasins{0}
  \def\HAFOCaputoShowReturn{0}
  \def\HAFOCaputoShowProtocol{0}
  \ifdefined\HAFOCaputoBlockMemory
    \def\HAFOCaputoShowMemory{1}
  \fi
  \ifdefined\HAFOCaputoBlockBasins
    \def\HAFOCaputoShowBasins{1}
  \fi
  \ifdefined\HAFOCaputoBlockReturn
    \def\HAFOCaputoShowReturn{1}
  \fi
  \ifdefined\HAFOCaputoBlockProtocol
    \def\HAFOCaputoShowProtocol{1}
  \fi
\fi

\ifnum\HAFOCaputoShowMemory=1\relax
\section{Geometría y topología de sistemas de Caputo: del punto a la memoria}
\label{sec:geometria-topologia-caputo}

\HAFOCardMemory

En los sistemas de Caputo, la memoria obliga a formular la dinámica en un
espacio funcional. La diferencia con las ecuaciones diferenciales ordinarias
se resume así:

\begin{center}
\fbox{\parbox{0.91\linewidth}{
En una ODE autónoma, el punto actual \(x(t)\in\mathbb R^n\) determina el
futuro. En un PVI de Caputo con terminal inferior finito, el punto actual no
contiene por sí solo la contribución acumulada del pasado. El objeto que
evoluciona como una semidinámica es un perfil funcional de memoria.
}}
\end{center}

La distinción no es terminológica. Cambia qué significa órbita, sección de
Poincaré, cuenca, invariancia, atractor, conjugación y punto perpetuo. Los
gráficos tridimensionales de los casos estudiados representan una
\emph{proyección} de una evolución funcional, no necesariamente el retrato de
fase completo.

\begin{hallazgo}{Memoria, atracción y espacio de fases}
La formulación de Volterra permite construir una semidinámica funcional,
aunque el estado puntual no genere en general una dinámica autónoma. Los
PVI con reinicio tampoco admiten soluciones no constantes exactamente
periódicas bajo las hipótesis de no periodicidad
\cite{CongTuan2017Nonlocal,DoanKloeden2021Semidynamical,
TavazoeiHaeri2009,KaslikSivasundaram2012}. La separación entre cuenca
funcional y cuenca de reinicio exige especificar las familias que se atraen:
el espacio completo de términos libres no admite atracción compacta de un
abierto. Para el Lorenz generalizado se demuestra aquí existencia global,
absorción y precompacidad de las órbitas de reinicio. Identificar sus
destinos, demostrar su atracción o establecer caos requiere argumentos
adicionales.
\end{hallazgo}

\subsection{Flujo ordinario en el espacio de estados}

Considérese primero la ODE autónoma
\begin{equation}
 \dot x=f(x),\qquad x(0)=x_0,\qquad f:\Omega\subset\mathbb R^n\to\mathbb R^n.
 \label{eq:geofo-ode}
\end{equation}
Si \(f\) es localmente Lipschitz, existe una solución local única. Cuando todas
las soluciones relevantes son globales se define
\begin{equation}
 \varphi_t(x_0):=x(t;x_0).
 \label{eq:geofo-ode-flow}
\end{equation}
La unicidad y la invariancia por traslación temporal implican
\begin{equation}
 \varphi_0=\operatorname{id},\qquad
 \varphi_{t+s}=\varphi_t\circ\varphi_s.
 \label{eq:geofo-flow-law}
\end{equation}
Si además se dispone de existencia única hacia tiempos negativos,
\(\{\varphi_t\}_{t\in\mathbb R}\) es un flujo; cada \(\varphi_t\) tiene inversa
\(\varphi_{-t}\). En ese marco, una órbita es
\(\{\varphi_t(x_0):t\in\mathbb R\}\), una sección transversal se corta con esa
órbita y las variedades estable e inestable se construyen usando los límites
hacia \(+\infty\) y \(-\infty\), respectivamente.

La ley \cref{eq:geofo-flow-law} expresa que detener la solución en \(s\) y
volver a arrancar desde \(x(s)\) no cambia el experimento. Ésta es precisamente
la propiedad que se pierde al conservar sólo el punto actual en Caputo.

\subsection{El PVI de Caputo y su ecuación de Volterra}

Para \(0<q<1\), considérese el sistema conmensurable
\begin{equation}
 {}^{\mathrm C}D_{0+}^{q}x(t)=f(x(t)),
 \qquad x(0)=x_0.
 \label{eq:geofo-caputo-pvi}
\end{equation}
Se ha tomado \(t_0=0\) sólo para aligerar la notación; reemplazar \(t\) por
\(t-t_0\) recupera un terminal inferior arbitrario. Bajo las hipótesis usuales,
\cref{eq:geofo-caputo-pvi} equivale a
\begin{equation}
 x(t)=x_0+\frac1{\Gamma(q)}
 \int_0^t(t-s)^{q-1}f(x(s))\,\mathrm ds.
 \label{eq:geofo-volterra}
\end{equation}
La integral contiene todos los valores \(f(x(s))\) entre \(0\) y \(t\). El
kernel \( (t-s)^{q-1}\) vuelve a ponderar esos valores cada vez que cambia el
tiempo de observación. Por ello no basta guardar una suma constante del pasado.

\begin{center}
\resizebox{0.96\linewidth}{!}{%
\begin{tikzpicture}[
  font=\small,
  >=Latex,
  caja/.style={draw=GrisOscuro,rounded corners=2pt,fill=Gris,inner sep=5pt,align=center},
  flecha/.style={-{Latex[length=2mm]},GrisOscuro}
]
  % Panel temporal: todos los tiempos pasados alimentan la observacion actual.
  \node[caja,anchor=west] at (0.15,4.35)
    {Forma de Volterra\\[-0.1em]
     \(\displaystyle x(t)=x_0+\int_0^t K_q(t-s)f(x(s))\,\mathrm ds\)};
  \draw[very thick,GrisOscuro] (0.65,0.95)--(7.15,0.95);
  \foreach \x/\lab in {0.65/{0},2.05/{s_1},3.45/{s_2},4.85/{s_3},6.15/{s_4},7.15/{t}}
    {\draw[GrisOscuro] (\x,0.82)--(\x,1.08) node[below=3pt] {\(\lab\)};}
  \node[circle,draw=Azul,very thick,fill=AzulClaro,minimum size=8mm] (obs) at (7.15,3.15) {\(x(t)\)};
  \draw[flecha,line width=0.45pt] (0.85,1.10) to[bend left=18] (obs.west);
  \draw[flecha,line width=0.60pt] (2.15,1.10) to[bend left=22] (obs.west);
  \draw[flecha,line width=0.80pt] (3.50,1.10) to[bend left=26] (obs.west);
  \draw[flecha,line width=1.05pt] (4.90,1.10) to[bend left=31] (obs.west);
  \draw[flecha,line width=1.35pt] (6.15,1.10) to[bend left=39] (obs.south);
  \node[align=center,GrisOscuro] at (3.55,0.12)
    {todo el intervalo \([0,t]\) contribuye; el grosor sólo recuerda\\
     que, para el kernel mostrado, los retardos pequeños tienen mayor peso};

  % Panel del kernel en funcion del retardo r=t-s.
  \begin{scope}[shift={(9.05,0.55)}]
    \draw[-{Latex[length=2mm]},GrisOscuro] (0,0)--(5.35,0)
      node[below left=1pt] {retardo \(r=t-s\)};
    \draw[-{Latex[length=2mm]},GrisOscuro] (0,0)--(0,3.75)
      node[above,align=center] {\(K_q(r)\)};
    \draw[Azul,very thick,domain=0.12:4.85,samples=90,smooth,variable=\r]
      plot ({\r},{1.18*pow(\r,-0.30)});
    \node[Azul,anchor=west] at (2.15,2.25) {\(q=0.7\) (forma normalizada)};
    \draw[dashed,GrisOscuro] (0.28,0)--(0.28,2.02);
    \node[align=left,anchor=west] at (0.45,3.25)
      {singularidad integrable\\cuando \(r\downarrow0\)};
    \draw[-{Latex[length=1.7mm]},GrisOscuro] (1.50,3.12)--(0.38,2.23);
    \node[align=left,anchor=west] at (2.55,0.75)
      {cola algebraica\\cuando crece \(r\)};
    \draw[-{Latex[length=1.7mm]},GrisOscuro] (3.58,1.15)--(4.45,1.02);
    \node[caja,anchor=north west] at (0.35,4.48)
      {\(\displaystyle K_q(r)=\frac{r^{q-1}}{\Gamma(q)},\quad 0<q<1\)};
  \end{scope}
\end{tikzpicture}%
}
\par\smallskip
\begin{minipage}{0.92\linewidth}
\small\textbf{Lectura conceptual.}
El kernel de Caputo no selecciona únicamente el último estado: vuelve a
ponderar el intervalo completo desde el terminal inferior. La curva usa
\(q=0.7\); no representa datos de ninguno de los casos de estudio.
\end{minipage}
\end{center}


Para garantizar existencia global y continuidad de la semidinámica funcional
es suficiente suponer que \(f\) es globalmente Lipschitz. En los sistemas
polinomiales analizados, que sólo son localmente Lipschitz, debe
añadirse una cota a priori que descarte explosión en tiempo finito o restringir
la construcción al intervalo máximo de existencia. Una trayectoria numérica
acotada en una ventana finita no sustituye esa prueba.

\subsection{Derivación del estado funcional de memoria}
\label{subsec:geofo-memory-profile}

Fijemos un tiempo \(t>0\) y escribamos el futuro como \(t+\theta\),
\(\theta\geq0\). Al separar la integral de \cref{eq:geofo-volterra} en
\([0,t]\) y \([t,t+\theta]\), se obtiene
\begin{align}
 x(t+\theta)
 ={}&x_0+\frac1{\Gamma(q)}
 \int_0^t(t+\theta-s)^{q-1}f(x(s))\,\mathrm ds
 \nonumber\\
 &+\frac1{\Gamma(q)}
 \int_t^{t+\theta}(t+\theta-s)^{q-1}f(x(s))\,\mathrm ds.
 \label{eq:geofo-split-memory}
\end{align}
Definimos el \emph{perfil de memoria acumulada}
\begin{equation}
 m_t(\theta):=x_0+\frac1{\Gamma(q)}
 \int_0^t(t+\theta-s)^{q-1}f(x(s))\,\mathrm ds,
 \qquad \theta\geq0.
 \label{eq:geofo-memory-profile}
\end{equation}
Este perfil no adivina el futuro: para cada retraso futuro \(\theta\) registra
cuál será la contribución del pasado ya conocido si se continúa hasta
\(t+\theta\). Haciendo \(s=t+u\) en la segunda integral de
\cref{eq:geofo-split-memory}, la continuación \(y_t(\theta):=x(t+\theta)\)
satisface
\begin{equation}
 y_t(\theta)=m_t(\theta)+\frac1{\Gamma(q)}
 \int_0^\theta(\theta-u)^{q-1}f(y_t(u))\,\mathrm du.
 \label{eq:geofo-shifted-volterra}
\end{equation}
En particular,
\begin{equation}
 m_t(0)=x_0+\frac1{\Gamma(q)}
 \int_0^t(t-s)^{q-1}f(x(s))\,\mathrm ds=x(t).
 \label{eq:geofo-evaluation-state}
\end{equation}
El punto \(x(t)\) es, pues, solamente el valor en cero de una función
\(m_t:[0,\infty)\to\mathbb R^n\). Conocer \(m_t(0)\) no determina en general
los valores \(m_t(\theta)\) para \(\theta>0\).

\begin{proposition}[el punto actual es una proyección]
Sea \(\operatorname{ev}_0(m)=m(0)\). La continuación exacta desde \(t\) está
determinada por \(m_t\), mientras que reiniciar el PVI con dato
\(x(t)=\operatorname{ev}_0(m_t)\) sustituye \(m_t\) por el perfil constante
\(\theta\mapsto x(t)\). Salvo casos especiales, ambas continuaciones son
distintas.
\end{proposition}

\begin{proof}
La continuación heredada resuelve \cref{eq:geofo-shifted-volterra}. El PVI
reiniciado en \(t\) resuelve
\begin{equation}
 z(\theta)=x(t)+\frac1{\Gamma(q)}
 \int_0^\theta(\theta-u)^{q-1}f(z(u))\,\mathrm du.
 \label{eq:geofo-reset-volterra}
\end{equation}
Los términos libres son \(m_t(\theta)\) y \(x(t)\), respectivamente. Sólo
coinciden si \(m_t\) es constante en la región relevante o si una cancelación
especial produce la misma solución.
\end{proof}

La posibilidad de que dos soluciones de un sistema Caputo de dimensión mayor
que uno alcancen el mismo punto con memorias diferentes es coherente con los
resultados de Cong y Tuan; en general no se genera una dinámica autónoma sobre
\(\mathbb R^n\) \cite{CongTuan2017Nonlocal}.

\begin{center}
\resizebox{0.96\linewidth}{!}{%
\begin{tikzpicture}[
  font=\small,
  >=Latex,
  perfil/.style={draw=GrisOscuro,rounded corners=3pt,minimum width=3.6cm,minimum height=1.15cm,align=center},
  futuro/.style={draw=GrisOscuro,rounded corners=3pt,minimum width=3.7cm,minimum height=1.15cm,align=center},
  arr/.style={-{Latex[length=2mm]},thick,GrisOscuro}
]
  \path[draw=GrisOscuro,rounded corners=4pt]
    (0.15,0.25) rectangle (4.45,4.85);
  \node[anchor=north west,fill=white,inner sep=2pt] at (0.38,4.68)
    {\(\mathfrak C\): perfiles funcionales};
  \node[perfil,fill=AzulClaro] (mone) at (2.30,3.55)
    {\(m_1\)\\perfil completo, trazo continuo};
  \node[perfil,fill=NaranjaClaro,dashed] (mtwo) at (2.30,1.55)
    {\(m_2\)\\perfil completo, trazo discontinuo};

  \node[circle,draw=GrisOscuro,very thick,fill=Gris,minimum size=1.25cm,align=center]
    (xstar) at (7.25,2.55) {\(x_\ast\)};
  \node[align=center] at (7.25,4.52) {proyección puntual\\en \(\mathbb R^n\)};

  \node[futuro,fill=AzulClaro] (yone) at (12.35,3.55)
    {continuación \(y_{m_1}\)\\determinada por \(m_1\)};
  \node[futuro,fill=NaranjaClaro,dashed] (ytwo) at (12.35,1.55)
    {continuación \(y_{m_2}\)\\determinada por \(m_2\)};
  \node[draw=Rojo,rounded corners=2pt,align=center,text=Rojo] (unknown) at (12.35,2.55)
    {el punto solo\\no selecciona un futuro};

  \draw[arr,Azul] (mone)--node[above,pos=0.56]{Volterra con \(m_1\)}(yone);
  \draw[arr,Naranja,dashed] (mtwo)--node[below,pos=0.56]{Volterra con \(m_2\)}(ytwo);
  \draw[arr] (mone.east) to[bend right=15] node[above,sloped]{\(\operatorname{ev}_0\)} (xstar.north west);
  \draw[arr,dashed] (mtwo.east) to[bend left=15] node[below,sloped]{\(\operatorname{ev}_0\)} (xstar.south west);
  \draw[-{Latex[length=2mm]},dashed,Rojo] (xstar)--(unknown);

  \node[draw=GrisOscuro,rounded corners=2pt,fill=Gris,align=center] at (7.25,0.22)
    {\(m_1(0)=m_2(0)=x_\ast\), pero \(m_1\neq m_2\) en \(\mathfrak C\)};
\end{tikzpicture}%
}
\par\smallskip
\begin{minipage}{0.92\linewidth}
\small\textbf{Estado frente a historia.}
En la realización funcional de Volterra pueden existir perfiles distintos con
la misma evaluación en cero. El dibujo expresa una posibilidad estructural;
no afirma que dos trayectorias concretas del banco numérico ya hayan exhibido
esa coincidencia.
\end{minipage}
\end{center}


\subsection{Contraejemplo explícito: Mittag--Leffler no es exponencial}

\ifHAFOPedagogy
\begin{bloqueestudio}{Ejercicios}
Compare numéricamente
\(E_q[-a(t+s)^q]\) y \(E_q(-at^q)E_q(-as^q)\) para \(q=0.8\).
Repita con \(q=1\) e interprete la diferencia.
\end{bloqueestudio}
\fi

El fallo de la ley de semigrupo puntual ya aparece en la ecuación escalar
\begin{equation}
 {}^{\mathrm C}D_{0+}^{q}x=\lambda x,\qquad x(0)=x_0.
 \label{eq:geofo-scalar-linear}
\end{equation}
Su solución es
\begin{equation}
 \Phi_t(x_0)=x_0E_q(\lambda t^q),
 \qquad
 E_q(z)=\sum_{k=0}^{\infty}\frac{z^k}{\Gamma(qk+1)}.
 \label{eq:geofo-mittag-solution}
\end{equation}
Si las aplicaciones puntuales formaran un semigrupo, sería necesario que
\begin{equation}
 E_q(\lambda(t+s)^q)
 =E_q(\lambda t^q)E_q(\lambda s^q).
 \label{eq:geofo-false-ml-product}
\end{equation}
Sin embargo, al comparar el término lineal en \(\lambda\), el miembro izquierdo
contiene
\((t+s)^q/\Gamma(q+1)\), mientras el derecho contiene
\((t^q+s^q)/\Gamma(q+1)\). Para \(0<q<1\) y \(t,s>0\), esas cantidades no son
iguales. La identidad reaparece cuando \(q=1\), pues \(E_1(z)=e^z\).

Este ejemplo separa dos operaciones:
\begin{align}
 \text{continuar: }&x(t+s)=
 \operatorname{ev}_0\bigl(S_sS_t\iota(x_0)\bigr),
 \label{eq:geofo-continue-operation}\\
 \text{reiniciar: }&z(s)=
 \operatorname{ev}_0\bigl(S_s\iota(
 \operatorname{ev}_0(S_t\iota(x_0)))\bigr),
 \label{eq:geofo-reset-operation}
\end{align}
donde \(S_t\) será la semidinámica funcional e
\(\iota(x)(\theta)\equiv x\) la inmersión de perfiles constantes. La segunda
expresión inserta \(\iota\circ\operatorname{ev}_0\): borra la memoria.

\subsection{Espacio de fases funcional y semidinámica}

\ifHAFOPedagogy
\begin{bloqueestudio}{Ejercicios}
Elija dos perfiles que coincidan en \([0,1]\) y difieran después de
\(\theta=2\); estime \(d_{\mathfrak C}\). Repita moviendo la diferencia a
\(\theta=20\).
\end{bloqueestudio}
\fi

Sea
\begin{equation}
 \mathfrak C:=C([0,\infty),\mathbb R^n)
 \label{eq:geofo-history-space}
\end{equation}
con la topología de convergencia uniforme en compactos. Una métrica que induce
esa topología es
\begin{equation}
 d_{\mathfrak C}(m,n)=
 \sum_{k=1}^{\infty}2^{-k}
 \frac{\displaystyle\sup_{0\leq\theta\leq k}
 \|m(\theta)-n(\theta)\|_2}
 {1+\displaystyle\sup_{0\leq\theta\leq k}
 \|m(\theta)-n(\theta)\|_2}.
 \label{eq:geofo-compact-open-metric}
\end{equation}
El uso de compactos significa que dos memorias son próximas si lo son en toda
ventana futura acotada; no se exige una cota uniforme sobre
\([0,\infty)\).

Para un perfil general \(m\in\mathfrak C\), considérese la ecuación de Volterra
\begin{equation}
 x_m(t)=m(t)+\frac1{\Gamma(q)}
 \int_0^t(t-s)^{q-1}f(x_m(s))\,\mathrm ds.
 \label{eq:geofo-general-volterra}
\end{equation}
Se define, para \(t,\theta\geq0\),
\begin{equation}
 (S_tm)(\theta)=m(t+\theta)+\frac1{\Gamma(q)}
 \int_0^t(t+\theta-s)^{q-1}f(x_m(s))\,\mathrm ds.
 \label{eq:geofo-semiflow-operator}
\end{equation}
Para \(m=\iota(x_0)\), esta fórmula coincide con
\cref{eq:geofo-memory-profile}. Además,
\begin{equation}
 \operatorname{ev}_0(S_tm)=x_m(t).
 \label{eq:geofo-evaluation-semiflow}
\end{equation}

\begin{definition}[semidinámica o semiflujo]
Una familia continua \(S:\mathbb R_+\times X\to X\) es una semidinámica si
\begin{equation}
 S_0=\operatorname{id}_X,
 \qquad S_{t+s}=S_t\circ S_s,
 \qquad t,s\geq0.
 \label{eq:geofo-semiflow-law}
\end{equation}
No se exige que \(S_t\) sea invertible. Si los tiempos negativos existen y
cada \(S_t\) es un homeomorfismo, se habla de flujo.
\end{definition}

\begin{proposition}[semidinámica Caputo en \(\mathfrak C\)]
Si \(f\) es globalmente Lipschitz, los operadores de
\cref{eq:geofo-semiflow-operator} forman una semidinámica continua sobre
\(\mathfrak C\) \cite{DoanKloeden2021Semidynamical}.
\end{proposition}

\begin{proof}[Idea algebraica de la propiedad de semigrupo]
Separar la integral en el tiempo \(t\) muestra primero que la curva
\(v\mapsto x_m(t+v)\) resuelve la ecuación de Volterra con término libre
\(S_tm\). La unicidad implica
\(x_{S_tm}(v)=x_m(t+v)\); ésta es la identidad que permite componer los
operadores.
Al calcular \(S_s(S_tm)\), la contribución heredada se separa en una integral
sobre \([0,t]\) y otra sobre \([t,t+s]\). En la segunda se hace el cambio
\(u=t+v\). Ambas integrales se unen en
\begin{equation*}
 \frac1{\Gamma(q)}\int_0^{t+s}
 (t+s+\theta-u)^{q-1}f(x_m(u))\,\mathrm du,
\end{equation*}
mientras el término libre es \(m(t+s+\theta)\). El resultado es
\((S_{t+s}m)(\theta)\). Si \(L\) es una constante Lipschitz de \(f\),
la dependencia respecto del término libre satisface, en cada \([0,T]\),
\[
 \sup_{0\leq t\leq T}\|x_m(t)-x_n(t)\|
 \leq\|m-n\|_{C([0,T])}E_q(LT^q),
\]
por la desigualdad de Grönwall fraccionaria. Junto con la continuidad de las
integrales de núcleo débilmente singular, esta cota proporciona la
continuidad de \(S\) en la topología compacto-abierta. Para la existencia
global, \(\|f(u)\|\le\|f(0)\|+L\|u\|\) y la misma desigualdad dan,
en cada horizonte finito \(T\),
\[
 \sup_{0\le t\le T}\|x_m(t)\|
 \le\left(\|m\|_{C([0,T])}
       +\frac{\|f(0)\|T^q}{\Gamma(q+1)}\right)E_q(LT^q).
\]
No puede ocurrir explosión en tiempo finito; el teorema de continuación de
Volterra prolonga la solución para todo \(t\ge0\)
\cite{DoanKloeden2021Semidynamical}.
\end{proof}

Esta ampliación tiene el siguiente alcance. Los perfiles constantes
\(\iota(x_0)\) representan exactamente los PVI convencionales de
\cref{eq:geofo-caputo-pvi}; un perfil general \(m\in\mathfrak C\) es un dato
admisible de la extensión de Volterra
\cref{eq:geofo-general-volterra}, pero no tiene por qué provenir de una
prehistoria física única. En consecuencia, toda afirmación de cuenca funcional
es relativa a esta realización semidinámica declarada. Además, restaurar el
semigrupo en todo \(\mathfrak C\) no proporciona automáticamente un espacio
adecuado para la atracción: los términos libres arbitrarios producen las
obstrucciones demostradas en \cref{prop:geofo-full-space-obstruction}.
La restauración de la propiedad de semigrupo no convierte todos los términos
libres continuos en memorias realizables por el sistema físico.

La geometría y la diferencia entre evolución funcional y reinicio puntual se
resumen en el diagrama siguiente.

\begin{center}
\resizebox{0.96\linewidth}{!}{%
\begin{tikzpicture}[
  font=\small,
  >=Latex,
  nodo/.style={draw=GrisOscuro,rounded corners=3pt,minimum width=2.55cm,minimum height=0.95cm,align=center},
  arr/.style={-{Latex[length=2mm]},thick,GrisOscuro}
]
  \node[draw=Verde,rounded corners=5pt,fill=VerdeClaro,fit={(0.15,3.05) (14.65,5.75)},
        label={[Verde,anchor=north west]north west:representación funcional en \(\mathfrak C\)}] {};
  \node[nodo,fill=white] (mzero) at (1.65,4.25) {\(m_0=\iota(x_0)\)};
  \node[nodo,fill=white] (mt) at (5.80,4.25) {\(m_t=S_tm_0\)};
  \node[nodo,fill=white] (mts) at (12.20,4.25) {\(m_{t+s}=S_s m_t\)};
  \draw[arr] (mzero)--node[above]{\(S_t\)}(mt);
  \draw[arr] (mt)--node[above]{\(S_s\)}(mts);
  \draw[arr,Verde] (mzero.north) to[bend left=24]
    node[above]{\(S_{t+s}=S_s\circ S_t\)} (mts.north);
  \node[draw=Verde,circle,very thick,fill=white] at (13.75,4.25) {\(\checkmark\)};

  \node[draw=Rojo,rounded corners=5pt,fill=NaranjaClaro,fit={(0.15,0.15) (14.65,2.85)},
        label={[Rojo,anchor=north west]north west:representación puntual en \(\mathbb R^n\)}] {};
  \node[nodo,fill=white] (xzero) at (1.65,1.20) {\(x_0\)};
  \node[nodo,fill=white] (xt) at (5.80,1.20)
    {\(x(t)=\operatorname{ev}_0m_t\)};
  \node[nodo,fill=white] (reset) at (9.55,1.75)
    {reinicio\\\(z(s)=\operatorname{ev}_0S_s\iota(x(t))\)};
  \node[nodo,fill=white] (inherit) at (12.20,0.55)
    {continuación\\\(x(t+s)=\operatorname{ev}_0m_{t+s}\)};

  \draw[arr] (xzero)--node[above]{\(\Phi_t=\operatorname{ev}_0S_t\iota\)}(xt);
  \draw[arr,dashed,Rojo] (xt)--node[above,sloped]{borra memoria}(reset);
  \draw[arr,Azul] (xzero.south) to[bend right=17]
    node[below]{\(\Phi_{t+s}\)} (inherit.west);
  \node[text=Rojo,font=\large] at (12.80,1.72) {\(\neq\) en general};

  \draw[arr] (xzero.north)--node[left]{\(\iota\)}(mzero.south);
  \draw[arr] (mt.south)--node[left]{\(\operatorname{ev}_0\)}(xt.north);
  \draw[arr] (mts.south)--node[right]{\(\operatorname{ev}_0\)}(inherit.north);
  \node[align=center,GrisOscuro] at (7.35,2.55)
    {la retracción \(\iota\circ\operatorname{ev}_0\) no es la identidad en \(\mathfrak C\)};
\end{tikzpicture}%
}
\par\smallskip
\begin{minipage}{0.92\linewidth}
\small\textbf{Dónde se restaura la composición.}
Bajo las hipótesis de existencia y unicidad declaradas, \(S_t\) compone sobre
perfiles. La familia puntual \(\Phi_t=\operatorname{ev}_0S_t\iota\) no forma,
en general, un semigrupo porque cada composición puntual introduce un reinicio.
\end{minipage}
\end{center}


\begin{table}[ht]
\centering
\small
\caption{Diccionario geométrico mínimo entre ODE y Caputo.}
\label{tab:geofo-dictionary}
\begin{tabularx}{\linewidth}{@{}L{0.25\linewidth}L{0.31\linewidth}Y@{}}
\toprule
Concepto & ODE autónoma & Caputo con terminal finito \\
\midrule
Estado completo & \(x\in\mathbb R^n\) & Perfil
\(m\in\mathfrak C\); \(x=\operatorname{ev}_0m\) es observable.\\
Evolución & Flujo o semiflujo \(\varphi_t\) puntual & Semidinámica \(S_t\)
funcional; el reinicio puntual no es composición de \(S_t\).\\
Órbita & Curva en \(\mathbb R^n\) & Curva \(t\mapsto m_t\) en
\(\mathfrak C\), proyectada como \(t\mapsto x(t)\).\\
Equilibrio & Punto \(e\) con \(f(e)=0\) & Perfil fijo constante
\(E_e(\theta)\equiv e\); en todo \(\mathfrak C\) hay además puntos fijos
artificiales.\\
Cuenca & Subconjunto del espacio de estados & Cuenca de reinicio o cuenca
relativa a un espacio admisible; la atracción abierta en todo
\(\mathfrak C\) está obstruida.\\
Poincaré & Mapa sobre una sección finito-dimensional & Operador sobre
\(\widehat\Sigma\subset\mathfrak C\); su proyección puede ser multivaluada.\\
Periodicidad & Ciclos permitidos & El PVI con reinicio no tiene soluciones
no constantes exactamente periódicas bajo las hipótesis finitas usuales.\\
\bottomrule
\end{tabularx}
\end{table}

\fi

\ifnum\HAFOCaputoShowBasins=1\relax
\ifdefined\HAFOCaputoBlockSelect
\section{Órbitas, invariancia y cuencas de sistemas de Caputo}
\fi
\subsection{Órbita funcional, proyección y conjunto límite}

La órbita positiva completa de \(m\in\mathfrak C\) es
\begin{equation}
 \mathcal O^+(m)=\{S_tm:t\geq0\}\subset\mathfrak C.
 \label{eq:geofo-functional-orbit}
\end{equation}
El retrato de fase observable es su proyección
\begin{equation}
 \mathcal O_x^+(m)=
 \{\operatorname{ev}_0(S_tm):t\geq0\}\subset\mathbb R^n.
 \label{eq:geofo-projected-orbit}
\end{equation}
La proyección puede cruzarse a sí misma sin que la órbita funcional lo haga:
dos perfiles distintos pueden tener el mismo valor en cero. Por consiguiente,
un cruce en un retrato \(x_i\)--\(x_j\) no implica ni falta de unicidad ni una
intersección en el espacio de fase verdadero.

Cuando la órbita es precompacta, su conjunto \(\omega\)-límite se define por
\begin{equation}
 \omega(m)=\bigcap_{T\geq0}
 \overline{\bigcup_{t\geq T}S_tm}^{\,\mathfrak C}.
 \label{eq:geofo-omega-limit}
\end{equation}
El conjunto observado
\(\operatorname{ev}_0[\omega(m)]\subset\mathbb R^n\) es sólo su sombra
finito-dimensional. Una nube de puntos acotada puede sugerir esa sombra, pero
no demuestra precompacidad en \(\mathfrak C\).

\subsection{Invariancia, atractor y cuenca en el espacio funcional}

Sea \(X\) un espacio métrico y \(S_t\) una semidinámica.

\begin{definition}[invariancia]
Un conjunto \(A\subset X\) es positivamente invariante si
\(S_tA\subseteq A\) para todo \(t\geq0\), y estrictamente invariante si
\(S_tA=A\). En una semidinámica, la igualdad es una condición más fuerte: no
se deduce de la existencia de inversas, porque éstas pueden no existir.
\end{definition}

\begin{definition}[atractor local]
Un conjunto \(A\subset X\) es un atractor local si es compacto, estrictamente
invariante y existe una vecindad abierta \(U\supset A\) tal que
\begin{equation}
 \operatorname{dist}_X(S_tm,A):=
 \inf_{a\in A}d_X(S_tm,a)\longrightarrow0
 \quad(t\to\infty)
 \label{eq:geofo-attraction}
\end{equation}
para todo \(m\in U\). Su cuenca es
\begin{equation}
 \mathcal B_X(A)=\{m\in X:
 \operatorname{dist}_X(S_tm,A)\to0\}.
 \label{eq:geofo-functional-basin}
\end{equation}
\end{definition}

Esta definición abstracta no puede aplicarse a un abierto del espacio completo
\(\mathfrak C\) para la extensión de Volterra anterior. Antes de precisar
la obstrucción, observemos que, si \(e\) satisface
\(f(e)=0\), el perfil constante
\begin{equation}
 E_e:=\iota(e),\qquad E_e(\theta)=e,
 \label{eq:geofo-equilibrium-profile}
\end{equation}
es un punto fijo: \(S_tE_e=E_e\). Así, un equilibrio puntual se convierte en
un equilibrio de la semidinámica funcional.

\begin{proposition}[obstrucciones en el espacio de todos los términos libres]
\label{prop:geofo-full-space-obstruction}
Supóngase \(f\) globalmente Lipschitz y considérese \(S_t\) sobre
\(\mathfrak C\) con la topología compacto-abierta.
\begin{enumerate}
 \item Ningún compacto no vacío atrae todos los perfiles de un abierto no vacío.
 \item Para cada \(c\in\mathbb R^n\), el perfil
 \begin{equation}
  m_c(\theta)=c-\frac{\theta^q}{\Gamma(q+1)}f(c)
  \label{eq:geofo-artificial-fixed-profile}
 \end{equation}
 es fijo, incluso cuando \(f(c)\ne0\).
 \item Si \(A\ne\varnothing\) es positivamente invariante, su conjunto
 de atracción puntual \(\mathcal B_{\mathfrak C}(A)\) es denso en
 \(\mathfrak C\).
\end{enumerate}
\end{proposition}

\begin{proof}
Sea \(U\) un abierto que contiene a \(m\). Elíjase un entero \(N\) tal
que \(2^{-N}\) sea menor que el radio de una bola centrada en \(m\)
contenida en \(U\). Cualquier perfil que coincida con \(m\) en \([0,N]\)
pertenece a esa bola, pues la cola de la métrica
\cref{eq:geofo-compact-open-metric} es a lo sumo \(2^{-N}\).
Escojamos una curva continua \(y\) igual a \(x_m\) hasta \(N\) y con
\(\|y(t)\|\to\infty\). El término libre
\[
 \widetilde m(t)=y(t)-\frac1{\Gamma(q)}
 \int_0^t(t-s)^{q-1}f(y(s))\,\mathrm ds
\]
es continuo, coincide con \(m\) hasta \(N\) y su solución es \(y\), por
unicidad. Si \(A\) es compacto, \(\operatorname{ev}_0(A)\) es acotado.
Por tanto, la diferencia entre \(\operatorname{ev}_0(S_t\widetilde m)=y(t)\)
y cualquier punto de \(\operatorname{ev}_0(A)\) diverge. Si
\(r_t=\operatorname{dist}(y(t),\operatorname{ev}_0(A))\), el sumando
\(k=1\) de la métrica da
\[
 \operatorname{dist}_{\mathfrak C}(S_t\widetilde m,A)
 \geq\frac12\frac{r_t}{1+r_t}.
\]
Como \(r_t\to\infty\), el límite inferior de esa distancia es al menos
\(1/2\), lo que impide la atracción.

Para el segundo enunciado, sustituir \(x_{m_c}(t)\equiv c\) en la ecuación
de Volterra verifica la solución constante. La integral en
\cref{eq:geofo-semiflow-operator} vale
\(f(c)[(t+\theta)^q-\theta^q]/\Gamma(q+1)\), de donde
\(S_tm_c=m_c\). Estos perfiles no representan equilibrios del PVI con
reinicio salvo que \(f(c)=0\).

Finalmente, dados \(m,N\) y \(a\in A\), construyamos \(y=x_m\) hasta
\(N\), interpolémosla continuamente hasta \(a(0)\) en \(T=N+1\) y
continuemos con \(y(T+\theta)=x_a(\theta)\). Con el mismo término libre
\(\widetilde m=y-I^qf(y)\), la separación de la integral da
\[
 (S_T\widetilde m)(\theta)
 =x_a(\theta)-\frac1{\Gamma(q)}
   \int_0^\theta(\theta-u)^{q-1}f(x_a(u))\,\mathrm du
 =a(\theta).
\]
La invariancia positiva implica que la órbita queda en \(A\) desde \(T\).
Así, toda vecindad de \(m\) contiene un perfil de su cuenca.
\end{proof}

En particular, para \(f(x)=-x\), los perfiles
\(m_c(\theta)=c[1+\theta^q/\Gamma(q+1)]\) son fijos y se aproximan a
\(E_0\) cuando \(c\to0\). La estabilidad asintótica del reinicio escalar
no implica atracción de \(E_0\) en todo \(\mathfrak C\). Asimismo, si
hay equilibrios, la densidad de la cuenca de cualquier conjunto no vacío
positivamente invariante impide excluir una vecindad abierta de \(E_e\).

\begin{definition}[atracción relativa a familias de reinicio]
\label{def:geofo-reset-family-attraction}
El universo de familias de reinicio acotadas es
\begin{equation}
 \mathcal D=\{\iota(B):B\subset\mathbb R^n\text{ acotado y no vacío}\}.
 \label{eq:geofo-reset-universe}
\end{equation}
Para estudiar un candidato local se especifica una subfamilia
\(\mathcal D_A\subseteq\mathcal D\) de datos situados en su región
candidata de atracción. Un compacto estrictamente invariante \(A\)
atrae esa familia si todas las soluciones correspondientes son globales y
\begin{equation}
 \sup_{x_0\in B}
 \operatorname{dist}_{\mathfrak C}(S_t\iota(x_0),A)\longrightarrow0,
 \qquad \iota(B)\in\mathcal D_A.
 \label{eq:geofo-relative-reset-attraction}
\end{equation}
La elección \(\mathcal D_A=\mathcal D\) expresa atracción de todas las
familias acotadas de reinicios; no se exige esta propiedad global a un
candidato local.
\end{definition}

La condición uniforme \cref{eq:geofo-relative-reset-attraction} es una
propiedad por demostrar, distinta tanto de la atracción punto a punto como
de atraer un abierto de \(\mathfrak C\). Esta distinción se estudia también
en el preprint \cite{CuiKloedenXin2026CaputoDynamics}. Su aplicación a un
sistema concreto exige verificar las hipótesis de atracción. Otra posibilidad es
construir un espacio de perfiles admisibles
\(X_{\mathrm{adm}}\subset\mathfrak C\), positivamente invariante y con un
semiflujo continuo en la topología relativa. Deben definirse la admisibilidad,
el dominio de existencia y las propiedades del cierre que aseguran la
continuidad de la evolución.

La proyección \(A_x=\operatorname{ev}_0(A)\) de un compacto funcional es
compacta porque \(\operatorname{ev}_0\) es continua. No obstante, no debe
declararse automáticamente \emph{atractor invariante en \(\mathbb R^n\)}: no
existe, en general, una semidinámica puntual respecto de la cual comprobar esa
invariancia.

El conjunto estable sí tiene una definición directa,
\begin{equation}
 W^s(A)=\{m\in\mathfrak C:
 \operatorname{dist}_{\mathfrak C}(S_tm,A)\to0\}.
 \label{eq:geofo-stable-set}
\end{equation}
En cambio, una \emph{variedad inestable} no puede escribirse ingenuamente con
\(S_{-t}\). Debe definirse mediante órbitas completas que existan en el
espacio funcional o mediante una teoría específica de variedades invariantes
para ecuaciones de Volterra. Ésta es una diferencia estructural con los
diagramas de separatrices de una ODE.

\subsection{Estabilidad espectral no equivale a geometría global}

Alrededor del equilibrio \(e\), la perturbación \(\xi=x-e\) satisface, al
primer orden,
\begin{equation}
 {}^{\mathrm C}D_{0+}^{q}\xi=J_f(e)\xi.
 \label{eq:geofo-linearization}
\end{equation}
Para orden conmensurable, el criterio de Matignon exige
\begin{equation}
 |\arg\lambda_j|>\frac{q\pi}{2}
 \qquad\text{para todo }\lambda_j\in\sigma(J_f(e))
 \label{eq:geofo-matignon}
\end{equation}
como condición de estabilidad asintótica lineal
\cite{Matignon1996}. Bajo hipótesis adicionales, la estabilidad de
Mittag--Leffler se extiende al sistema no lineal
\cite{LiChenPodlubny2009,CongTuanTrinh2020Asymptotic}. Estas conclusiones
se refieren al PVI con reinicio; no implican estabilidad atractiva en el
espacio completo de términos libres, como demuestra
\cref{eq:geofo-artificial-fixed-profile}.

Geométricamente, \cref{eq:geofo-matignon} describe un sector angular del plano
espectral. No determina por sí sola:
\begin{itemize}
 \item la forma ni la conectividad de la cuenca;
 \item la existencia de una separatriz global;
 \item la coexistencia de atractores;
 \item la ocultedad de una cuenca;
 \item una herradura, entropía topológica o transitividad.
\end{itemize}
Además, el decaimiento típico es de Mittag--Leffler y puede ser algebraico, no
exponencial. Por eso la dicotomía hiperbólica y el lema de la variedad estable
de una ODE no se transfieren sustituyendo simplemente \(e^{At}\) por
\(E_q(At^q)\).

\subsection{Cuenca funcional frente a sección de reinicio}

\ifHAFOPedagogy
\begin{bloqueestudio}{Ejercicios}
Dibuje \(\iota(\mathbb R^3)\subset\mathfrak C\) para Chua. Compare una bola
euclídea alrededor de un equilibrio con una bola funcional y construya una
perturbación de historia que conserve el valor terminal. Explique por qué
esta perturbación no pertenece automáticamente a un espacio de perfiles
admisibles y por qué la densidad de la cuenca en todo \(\mathfrak C\) no
implica densidad de su sección de reinicio.
\end{bloqueestudio}
\fi
\label{subsec:geofo-two-basins}

La inmersión de reinicio
\begin{equation}
 \iota:\mathbb R^n\longrightarrow\mathfrak C,
 \qquad [\iota(x)](\theta)=x,
 \label{eq:geofo-reset-embedding}
\end{equation}
selecciona una subvariedad finito-dimensional de perfiles constantes. La
cuenca que realmente se muestrea al iniciar muchas simulaciones convencionales
es
\begin{equation}
 \mathcal B_{\mathrm{reset}}(A)
 :=\iota^{-1}\!\left[\mathcal B_{\mathfrak C}(A)\right]
 \subset\mathbb R^n.
 \label{eq:geofo-reset-basin}
\end{equation}
Aquí se supone que el semiflujo está definido en los perfiles considerados.
En un espacio restringido \(X_{\mathrm{adm}}\), se usa
\(\iota^{-1}[\mathcal B_{X_{\mathrm{adm}}}(A)]\) y se incluyen sólo los
reinicios cuya imagen pertenece a dicho espacio. La igualdad de conjuntos
expresa atracción puntual; no demuestra la atracción uniforme de familias
de \cref{eq:geofo-relative-reset-attraction}.

\begin{definition}[ocultedad funcional relativa a perfiles admisibles]
Supóngase construido un espacio \(X_{\mathrm{adm}}\subset\mathfrak C\)
positivamente invariante, con semiflujo continuo en la topología relativa y
que contenga los perfiles \(E_e\) de los equilibrios considerados.
Sea \(A\subset X_{\mathrm{adm}}\) un atractor local en ese espacio.
Diremos que \(A\) es funcionalmente oculto respecto de esos equilibrios si
para cada \(e\) existe una vecindad \(U_e\) de \(E_e\), abierta en
\(X_{\mathrm{adm}}\), tal que
\begin{equation}
 U_e\cap\mathcal B_{X_{\mathrm{adm}}}(A)=\varnothing.
 \label{eq:geofo-functional-hiddenness}
\end{equation}
\end{definition}

\begin{definition}[ocultedad sobre la sección de reinicio]
Sea \(A\) un compacto invariante cuya atracción respecto de una familia
\(\mathcal D_A\) esté establecida. Diremos que \(A\) es oculto bajo la
convención de reinicio si para cada
equilibrio \(e\) existe \(\rho_e>0\) tal que
\begin{equation}
 B_{\mathbb R^n}(e,\rho_e)
 \cap\mathcal B_{\mathrm{reset}}(A)=\varnothing.
 \label{eq:geofo-reset-hiddenness}
\end{equation}
\end{definition}

La primera condición implica la exclusión de cuenca de
\cref{eq:geofo-reset-hiddenness} si, para cada equilibrio, existe una
vecindad euclídea \(V_e\) tal que \(\iota(V_e)\subset X_{\mathrm{adm}}\)
y se usa la misma cuenca relativa. En efecto, la continuidad de
\(\iota:V_e\to X_{\mathrm{adm}}\) hace que la preimagen de \(U_e\)
contenga una bola alrededor de \(e\). La implicación inversa no se deduce:
los perfiles constantes no agotan las direcciones funcionales admisibles.
La atracción uniforme de \(\mathcal D_A\) sigue siendo una hipótesis
separada. Bajo esas condiciones,
\begin{equation}
 \text{ocultedad funcional}
 \Longrightarrow
 \text{ocultedad de reinicio},
 \qquad
 \text{pero no a la inversa.}
 \label{eq:geofo-hiddenness-hierarchy}
\end{equation}

\begin{center}
\resizebox{0.96\linewidth}{!}{%
\begin{tikzpicture}[font=\small,>=Latex]
  \path[draw=GrisOscuro,very thick,rounded corners=8pt,fill=Gris]
    (0.20,0.35) rectangle (14.70,5.45);
  \node[anchor=north west,font=\bfseries] at (0.45,5.25)
    {espacio funcional \(\mathfrak C\) (representación esquemática)};

  % Cuenca funcional: la lengua superior entra en U_e sin tocar el corte cercano.
  \path[draw=Verde,very thick,fill=VerdeClaro]
    (4.35,3.05)
    .. controls (5.05,3.85) and (6.15,3.90) .. (7.05,3.30)
    .. controls (8.15,4.70) and (11.85,4.80) .. (13.55,3.65)
    .. controls (14.15,3.20) and (14.05,1.70) .. (13.20,1.25)
    .. controls (11.90,0.65) and (9.10,0.90) .. (7.95,1.75)
    .. controls (7.25,2.25) and (7.00,2.65) .. (6.35,2.75)
    .. controls (5.55,2.88) and (4.85,2.55) .. (4.35,3.05) -- cycle;
  \node[Verde,align=center] at (10.70,3.55)
    {cuenca funcional\\\(\mathcal B_{\mathfrak C}(A)\)};
  \node[draw=Verde,very thick,circle,fill=white,minimum size=8mm] (A) at (12.35,2.55) {\(A\)};

  % Corte finito-dimensional de perfiles constantes.
  \draw[GrisOscuro,very thick]
    (0.85,1.55) .. controls (4.15,1.95) and (8.40,1.35) .. (13.95,1.75);
  \node[anchor=west,fill=white,inner sep=2pt] (slice) at (0.70,2.55)
    {sección de reinicio \(\iota(\mathbb R^n)\)};
  \draw[-{Latex[length=1.7mm]},GrisOscuro] (slice.south)--(1.65,1.66);
  \draw[Verde,line width=3pt]
    (8.30,1.50) .. controls (10.10,1.40) and (12.05,1.48) .. (13.55,1.69);
  \node[Verde,anchor=south,fill=white,inner sep=1pt] at (10.65,1.62)
    {\(\iota[\mathcal B_{\mathrm{reset}}(A)]\)};

  % Equilibrio y vecindad funcional.
  \node[draw=Azul,very thick,circle,fill=AzulClaro,minimum size=7mm] (Ee) at (3.70,1.82) {\(E_e\)};
  \draw[Azul,very thick,dashed] (3.70,1.82) ellipse (1.25cm and 1.42cm);
  \node[Azul,anchor=east,align=right] at (2.35,3.00)
    {vecindad funcional\\\(U_e\)};
  \draw[-{Latex[length=2mm]},Azul] (2.45,2.83)--(2.85,2.55);
  \node[Naranja,align=center] at (5.15,4.40)
    {perfil no constante cercano\\que pertenece a la cuenca};
  \draw[-{Latex[length=2mm]},Naranja,thick] (5.15,4.05)--(4.72,3.12);

  \node[draw=GrisOscuro,rounded corners=2pt,fill=white,align=center] at (7.45,0.60)
    {ocultedad funcional \(\Longrightarrow\) ocultedad de reinicio;
     el recíproco no vale en general};
\end{tikzpicture}%
}
\par\smallskip
\begin{minipage}{0.92\linewidth}
\small\textbf{Esquema, no cuenca calculada.}
La figura muestra una configuración compatible con la jerarquía lógica: una
bola del corte de reinicio puede evitar la cuenca mientras una perturbación
funcional cercana sí la alcanza. No atribuye esta geometría a PLL, MAVPD, Wu ni
a otro caso del banco actual.
\end{minipage}
\end{center}


\begin{proposalbox}
La ocultedad funcional en todo \(\mathfrak C\) está descartada por
\cref{prop:geofo-full-space-obstruction} cuando existen equilibrios; su
formulación relativa requiere construir y justificar
\(X_{\mathrm{adm}}\). Los sondeos finitos aproximan la condición
\cref{eq:geofo-reset-hiddenness} con radios, malla, integrador, terminal
inferior y política de memoria declarados. Esos sondeos no verifican toda una
bola de reinicios ni prueban por sí solos la atracción uniforme de
\(\mathcal D_A\). Los sondeos no constantes tampoco sustituyen la
construcción del espacio admisible.
\end{proposalbox}

\subsection{Estabilidad local y separación de cuencas de reinicio}

La estabilidad de los equilibrios puede aportar una exclusión de cuenca
demostrable. El paso necesario consiste en relacionar la convergencia de
estados con la de sus perfiles heredados.

\begin{proposition}[separación respecto de un equilibrio atractivo]
\label{prop:geofo-reset-equilibrium-separation}
Sea \(f\) continuo y supóngase bien planteado el PVI con reinicio. Si una
solución global satisface \(x(t)\to e\), con \(f(e)=0\), entonces
\(m_t\to E_e\) en \(\mathfrak C\). En consecuencia, si todos los
reinicios de una vecindad \(V_e\) convergen a \(e\) y un compacto
\(A\subset\mathfrak C\) no contiene a \(E_e\), se tiene
\[
 V_e\cap\mathcal B_{\mathrm{reset}}(A)=\varnothing.
\]
\end{proposition}

\begin{proof}
La ecuación de continuación permite escribir
\[
 m_t(\theta)-e=x(t+\theta)-e
 -\frac1{\Gamma(q)}\int_0^\theta
   (\theta-s)^{q-1}f(x(t+s))\,\mathrm ds.
\]
Por tanto, para cada \(\Theta>0\),
\begin{align*}
 \sup_{0\leq\theta\leq\Theta}\|m_t(\theta)-e\|
 &\leq\sup_{u\in[t,t+\Theta]}\|x(u)-e\|\\
 &\quad+\frac{\Theta^q}{\Gamma(q+1)}
       \sup_{u\in[t,t+\Theta]}\|f(x(u))\|
 \longrightarrow0.
\end{align*}
Esto es convergencia uniforme en compactos. Como \(A\) es compacto y
\(E_e\notin A\), la distancia
\(\operatorname{dist}_{\mathfrak C}(E_e,A)\) es positiva.
Una órbita de perfiles que converge a \(E_e\) no puede aproximarse también
a \(A\), lo que prueba la exclusión.
\end{proof}

Si todos los equilibrios satisfacen un teorema de estabilidad asintótica local
para reinicios, y un compacto invariante \(A\) no contiene ninguno de sus
perfiles, esta proposición establece la separación exigida por
\cref{eq:geofo-reset-hiddenness}. Para clasificar \(A\) como atractor
oculto queda por demostrar su atracción respecto de la familia
\(\mathcal D_A\). Así, Matignon y el teorema no lineal correspondiente
pueden excluir regiones próximas a los equilibrios, mientras la localización
del candidato debe buscarse fuera de ellas. La conclusión concierne a
reinicios; no establece ocultedad funcional en todo \(\mathfrak C\).

Si el sistema no tiene equilibrios, como el candidato de Muñoz--Pacheco, la
condición basada en vecindades de equilibrios se vuelve vacía. Es preferible
informar \emph{sistema sin equilibrios con conjunto atrayente observado} y
estudiar su cuenca funcional, en vez de interpretar la vacuidad lógica como
una prueba geométrica fuerte de ocultedad.

\fi

\ifnum\HAFOCaputoShowReturn=1\relax
\ifdefined\HAFOCaputoBlockSelect
\section{Retorno, recurrencia y puntos perpetuos con memoria}
\fi
\subsection{Cambios de coordenadas, conjugación y geometría diferencial}

En una ODE sobre una variedad diferenciable, la regla de la cadena hace que un
campo vectorial se transforme correctamente entre cartas. Para Caputo, en
general,
\begin{equation}
 {}^{\mathrm C}D^q[h(x(t))]
 \neq Dh(x(t))\,{}^{\mathrm C}D^qx(t).
 \label{eq:geofo-no-chain-rule}
\end{equation}
Por ello, una ecuación Caputo definida componente a componente en una carta no
produce automáticamente una ecuación del mismo tipo bajo un cambio no lineal
de coordenadas. El contraejemplo explícito de
\cref{eq:pp-chain-counterexample} muestra esta falla.

Los cambios afines sí son seguros. Si \(x=Tz+c\), con \(T\) constante e
invertible, entonces
\begin{equation}
 {}^{\mathrm C}D^qx=T\,{}^{\mathrm C}D^qz,
 \qquad
 {}^{\mathrm C}D^qz=T^{-1}f(Tz+c).
 \label{eq:geofo-affine-transform}
\end{equation}
El homeomorfismo funcional
\begin{equation}
 (H_{T,c}m)(\theta)=T^{-1}(m(\theta)-c)
 \label{eq:geofo-affine-history-map}
\end{equation}
conjuga las semidinámicas correspondientes:
\begin{equation}
 H_{T,c}\circ S_t=\widetilde S_t\circ H_{T,c}.
 \label{eq:geofo-affine-conjugacy}
\end{equation}
Esta igualdad fundamenta la transformación afín de atractores y cuencas
empleada en los análisis.

Abstractamente, dos semidinámicas \(S_t\) sobre \(X\) y \(\widetilde S_t\)
sobre \(Y\) son topológicamente conjugadas si existe un homeomorfismo
\(H:X\to Y\) que satisface
\begin{equation}
 H\circ S_t=\widetilde S_t\circ H,
 \qquad t\geq0.
 \label{eq:geofo-topological-conjugacy}
\end{equation}
La conjugación funcional transporta invariancia, recurrencia y propiedades
topológicas. Un homeomorfismo puntual
\(h:\mathbb R^n\to\mathbb R^n\) puede inducir una conjugación abstracta de
perfiles, pero la semidinámica transformada no tiene por qué escribirse como
una ecuación autónoma local \({}^{\mathrm C}D^qy=g(y)\).

\begin{remark}[problema geométrico abierto]
Una formulación intrínseca de derivadas fraccionarias sobre una variedad debe
especificar cómo comparar datos situados en espacios tangentes de tiempos
distintos, por ejemplo mediante una conexión, transporte paralelo o una
inmersión ambiente. El Caputo componente a componente es una formulación
extrínseca en \(\mathbb R^n\) y depende de la carta física elegida.
\end{remark}

\subsection{Secciones de Poincaré: de mapa puntual a operador funcional}

\ifHAFOPedagogy
\begin{bloqueestudio}{Ejercicios}
Guarde en cada cruce \(x(t_k)\) y un vector de memoria. Determine si dos cruces
cercanos en estado conservan sucesores cercanos al condicionar por la historia.
\end{bloqueestudio}
\fi

Para una ODE, una hipersuperficie
\begin{equation}
 \Sigma=\{x\in\mathbb R^n:h(x)=0\}
 \label{eq:geofo-point-section}
\end{equation}
es transversal cuando \(\nabla h(x)\cdot f(x)\neq0\). El primer retorno define
un mapa \(P:\Sigma\to\Sigma\). En Caputo, dos dificultades impiden copiar esta
construcción sin cambios:
\begin{enumerate}
 \item la continuación depende del perfil \(m\), no sólo de
 \(x=\operatorname{ev}_0m\);
 \item la velocidad ordinaria \(x'(t)\) no es \(f(x(t))\), y puede ser singular
 cerca del terminal inferior.
\end{enumerate}

La sección observable se levanta a
\begin{equation}
 \widehat\Sigma=
 \{m\in\mathfrak C:h(\operatorname{ev}_0m)=0\}.
 \label{eq:geofo-functional-section}
\end{equation}
Para un perfil \(m\) que tenga un siguiente cruce orientado, se define
\begin{equation}
 \tau(m)=\inf\{t>\tau_{\min}:
 h(\operatorname{ev}_0S_tm)=0
 \text{ con la orientación prescrita}\},
 \label{eq:geofo-return-time}
\end{equation}
y el operador de retorno funcional
\begin{equation}
 \mathcal P(m)=S_{\tau(m)}m,
 \qquad \mathcal P:\widehat\Sigma\dashrightarrow\widehat\Sigma.
 \label{eq:geofo-functional-poincare}
\end{equation}
La flecha discontinua recuerda que no todos los perfiles regresan y que
\(\tau\) puede perder continuidad cerca de tangencias.

En cómputo, la orientación debe detectarse por cambio de signo y refinamiento
del evento; no por la fórmula ODE \(\nabla h\cdot f\). Si \(m_1,m_2\) tienen
el mismo punto actual pero memorias distintas, es posible que
\begin{equation}
 \operatorname{ev}_0\mathcal P(m_1)\neq
 \operatorname{ev}_0\mathcal P(m_2).
 \label{eq:geofo-multivalued-projection}
\end{equation}
Así, los pares graficados en \(\Sigma\subset\mathbb R^n\) forman en general
una \emph{relación de retorno proyectada},
\begin{equation}
 \mathcal R_\Sigma=
 \{(\operatorname{ev}_0m,
 \operatorname{ev}_0\mathcal P(m)):m\in\widehat\Sigma\},
 \label{eq:geofo-projected-return-relation}
\end{equation}
no un mapa univaluado demostrado. Para recuperar un mapa finito-dimensional
sería necesario probar, por ejemplo, una reconstrucción única de memoria desde
las coordenadas observadas o la existencia de una variedad inercial adecuada.

\begin{center}
\resizebox{0.96\linewidth}{!}{%
\begin{tikzpicture}[
  font=\small,
  >=Latex,
  punto/.style={circle,draw=GrisOscuro,very thick,fill=white,minimum size=6mm},
  arr/.style={-{Latex[length=2mm]},thick,GrisOscuro}
]
  % Panel ODE: transversalidad diferencial clasica.
  \path[draw=GrisOscuro,rounded corners=5pt,fill=Gris]
    (0.15,0.25) rectangle (5.05,5.50);
  \node[anchor=north west,font=\bfseries] at (0.38,5.28) {ODE: sección transversal};
  \draw[Azul,very thick] (2.60,0.75)--(2.60,4.60) node[above] {\(\Sigma\)};
  \draw[GrisOscuro,very thick,-{Latex[length=2mm]}]
    (0.65,1.15) .. controls (1.45,2.00) and (1.70,3.05) .. (2.60,3.05)
    .. controls (3.45,3.02) and (4.00,3.65) .. (4.55,4.45);
  \node[punto] (cross) at (2.60,3.05) {};
  \draw[-{Latex[length=2mm]},very thick,Naranja] (cross)--(3.72,3.46)
    node[above right] {\(f(x)\)};
  \draw[-{Latex[length=2mm]},very thick,Verde] (cross)--(1.55,3.05)
    node[above] {\(\nabla h\)};
  \node[align=center] at (2.55,0.55)
    {\(\nabla h(x)\!\cdot\! f(x)\neq0\)};

  % Panel funcional: retornos distintos se colapsan al proyectar el punto inicial.
  \path[draw=GrisOscuro,rounded corners=5pt,fill=white]
    (5.35,0.25) rectangle (14.75,5.50);
  \node[anchor=north west,font=\bfseries] at (5.58,5.28)
    {Caputo: operador parcial sobre \(\widehat\Sigma\subset\mathfrak C\)};
  \path[draw=Azul,very thick,rounded corners=12pt,fill=AzulClaro]
    (5.85,2.30) rectangle (14.25,4.62);
  \node[Azul,anchor=north east] at (14.05,4.48) {\(\widehat\Sigma\)};
  \node[punto] (mone) at (7.15,4.05) {\(m_1\)};
  \node[punto,dashed] (mtwo) at (7.15,3.05) {\(m_2\)};
  \node[punto] (pone) at (12.65,4.12) {\(\mathcal Pm_1\)};
  \node[punto,dashed] (ptwo) at (12.65,2.98) {\(\mathcal Pm_2\)};
  \draw[arr] (mone)--node[above]{primer retorno, si existe}(pone);
  \draw[arr,dashed] (mtwo)--(ptwo);

  \draw[GrisOscuro,very thick] (6.05,0.90)--(14.05,0.90)
    node[right] {\(\Sigma\subset\mathbb R^n\)};
  \node[punto,fill=Gris] (xstar) at (7.15,0.90) {\(x_\ast\)};
  \node[punto] (xone) at (12.15,0.90) {\(x_1^+\)};
  \node[punto,dashed] (xtwo) at (13.30,0.90) {\(x_2^+\)};
  \draw[arr] (mone)--node[left]{\(\operatorname{ev}_0\)}(xstar);
  \draw[arr,dashed] (mtwo) to[bend left=14] (xstar);
  \draw[arr] (pone)--(xone);
  \draw[arr,dashed] (ptwo)--(xtwo);
  \node[align=center,fill=white,inner sep=2pt] at (9.55,1.72)
    {\(m_1(0)=m_2(0)=x_\ast\)\\pero \(x_1^+\) y \(x_2^+\) pueden diferir};
\end{tikzpicture}%
}
\par\smallskip
\begin{minipage}{0.92\linewidth}
\small\textbf{Transversalidad y retorno.}
En la ODE la orientación local se obtiene del campo tangente. En Caputo el
cruce observable se localiza causalmente por cambio de signo refinado y el
retorno conserva el perfil. El operador es parcial: algunos perfiles pueden no
regresar y la proyección sólo será univaluada si se demuestra una reducción
adicional de la memoria.
\end{minipage}
\end{center}


\subsection{Periodicidad exacta y recurrencia: distinciones terminológicas}

Para terminal inferior finito, orden no entero y las hipótesis regulares
consideradas en la literatura, una solución no constante del PVI con reinicio
de un sistema Caputo autónomo no puede ser exactamente periódica
\cite{TavazoeiHaeri2009,KaslikSivasundaram2012}. La razón estructural es que la
derivada fraccionaria de una función periódica no conserva, en general, el
mismo período; el kernel sigue recordando el arranque.
La afirmación no se extiende a soluciones con términos libres arbitrarios
\(m\in\mathfrak C\), que pueden construirse a partir de una curva
prescrita mediante \(m=y-I^qf(y)\).

Esto no impide observar curvas cerradas con precisión gráfica, oscilaciones
persistentes o recurrencia muy fuerte. Estas observaciones motivan dos medidas
distintas:
\begin{align}
 R_x(T;t)&=\|x(t+T)-x(t)\|_2,
 \label{eq:geofo-point-recurrence}\\
 R_{\mathfrak C}(T;t)&=
 d_{\mathfrak C}(m_{t+T},m_t).
 \label{eq:geofo-memory-recurrence}
\end{align}
Un \(R_x\) pequeño no obliga a que \(R_{\mathfrak C}\) sea pequeño. Según la
evidencia disponible, la terminología aplicable incluye \emph{oscilación acotada},
\emph{recurrencia proyectada}, \emph{lazo aparentemente periódico en una
ventana} o \emph{régimen asintóticamente periódico}, según la evidencia. La
expresión \emph{ciclo límite Caputo exacto} requiere cambiar el operador, llevar
el terminal inferior a un régimen bilateral apropiado o demostrar que se está
fuera de las hipótesis de no periodicidad.

Tampoco se deduce que el caos fraccionario no existe. Lo que cambia es el
objeto sobre el cual deben formularse sensibilidad, transitividad, entropía y
recurrencia. Una definición topológica rigurosa debe aplicarse a \(S_t\) o a un
operador funcional como \(\mathcal P\), no únicamente a la nube proyectada.

\subsection{Exponentes, entropía y límites de la evidencia numérica}

Para $f\in C^1$ en una región que contenga las trayectorias próximas
durante un horizonte finito, la derivación respecto del dato inicial de
la ecuación de Volterra da
\begin{equation}
 Y(t)=I+\frac1{\Gamma(q)}\int_0^t(t-s)^{q-1}J_f(x(s))Y(s)\,\mathrm ds.
 \label{eq:geofo-tangent-volterra}
\end{equation}
Para justificar el paso, se consideran cocientes incrementales del dato
inicial. El teorema fundamental del cálculo expresa la diferencia del
campo mediante la integral de $Df$ sobre el segmento entre ambas
soluciones. La continuidad uniforme de $Df$ en el compacto y Gronwall
fraccionario hacen converger esos cocientes a la única solución de
\cref{eq:geofo-tangent-volterra}. En campos no suaves esta diferenciabilidad
requiere un argumento adicional. La ecuación variacional a lo largo de
la trayectoria es
\begin{equation}
 {}^{\mathrm C}D_{0+}^{q}\xi(t)=J_f(x(t))\xi(t).
\label{eq:geofo-fractional-variational}
\end{equation}
El efecto de la reortogonalización se obtiene algebraicamente. Si
$Y(\tau)=QR$, con $R$ invertible, el cambio de base $Z(t)=Y(t)R^{-1}$
satisface
\[
 Z(t)=R^{-1}+\frac1{\Gamma(q)}\int_0^t(t-s)^{q-1}
 J_f(x(s))Y(s)R^{-1}\,\mathrm ds.
\]
Además del valor $Z(\tau)=Q$, deben transformarse el dato inicial tangente
y todos los valores históricos del integrando. Reiniciar en $\tau$ usando
sólo $Q$ elimina la contribución del pasado, que depende del tiempo futuro
cuando $0<q<1$. Para $q=1$ dicha contribución se reduce al estado tangente
en $\tau$. La invertibilidad de $R$ se comprueba: no se atribuyen sin
demostración a la matriz fundamental fraccionaria todas las propiedades
de la correspondiente matriz de una EDO.
También ella conserva memoria. Un exponente positivo calculado con una técnica
fraccionaria apropiada aporta evidencia de crecimiento de perturbaciones, pero
no demuestra por sí solo una herradura de Smale, conjugación a un corrimiento,
entropía topológica positiva ni ocultedad. Los algoritmos que reortogonalizan o
reinician vectores tangentes deben especificar cómo preservan la memoria de
\cref{eq:geofo-fractional-variational}; la metodología de dinámicas clonadas es
un ejemplo de tratamiento especializado \cite{fischer2020}.

De forma análoga:
\begin{itemize}
 \item una FFT de banda ancha no prueba caos topológico;
 \item una sección proyectada con muchos puntos no prueba un mapa de Poincaré
 finito-dimensional;
 \item una dimensión fractal estimada en \(\mathbb R^3\) caracteriza la
 proyección muestreada, no automáticamente el atractor en \(\mathfrak C\);
 \item un test \(0\)--\(1\) o un exponente de Lyapunov es diagnóstico dinámico,
 no certificado de cuenca.
\end{itemize}

\subsection{Puntos perpetuos vistos desde el espacio funcional}

Las dos construcciones de puntos perpetuos de \cref{sec:perpetual-points}
actúan sobre objetos geométricos diferentes.

\paragraph{FPP--A: semilla sobre la sección de reinicio.}
Una raíz
\begin{equation}
 J_f(p)f(p)=0,\qquad f(p)\neq0
 \label{eq:geofo-fppa}
\end{equation}
se inserta como perfil constante \(\iota(p)\). La condición cancela el
coeficiente \(h^{2q}\) de la expansión Picard--Caputo al arrancar desde ese
perfil. Por ello FPP--A describe geometría local de la salida desde la sección
de reinicio; no define una aceleración global ni un conjunto invariante de
\(S_t\).

\paragraph{FPP--H: evento sobre una trayectoria con memoria.}
El objeto
\begin{equation}
 \mathcal A_{q,q}[x](t)=
 {}^{\mathrm C}D_{0+}^{q}[f(x(\cdot))](t)
 \label{eq:geofo-fpph}
\end{equation}
depende del recorrido completo entre \(0\) y \(t\). Un cero con
\(f(x(t))\neq0\) es un evento de la pareja estado--historia, no una raíz en
\(\mathbb R^n\). Para evaluarlo numéricamente se conserva además de \(m_t\) la
serie histórica necesaria para la convolución L1. El perfil \(m_t\) es
suficiente para continuar el futuro, pero no se presupone que permita
reconstruir de manera única todos los valores pasados usados por el diagnóstico
secuencial.

\begin{center}
\resizebox{0.96\linewidth}{!}{%
\begin{tikzpicture}[
  font=\small,
  >=Latex,
  box/.style={draw=GrisOscuro,rounded corners=3pt,align=center,inner sep=5pt},
  arr/.style={-{Latex[length=2mm]},thick,GrisOscuro}
]
  % Carril algebraico FPP-A.
  \path[draw=Azul,very thick,rounded corners=6pt,fill=AzulClaro]
    (0.15,0.95) rectangle (7.15,7.00);
  \node[font=\bfseries,Azul] at (3.65,6.67)
    {FPP--A: semilla algebraica};
  \node[box,fill=white] (root) at (2.00,5.35)
    {\(J_f(p)f(p)=0\)\\\(f(p)\neq0\)};
  \node[box,fill=white] (embed) at (5.35,5.35)
    {perfil constante\\\(\iota(p)\in\mathfrak C\)};
  \draw[arr] (root)--node[above]{\(\iota\)}(embed);
  \node[box,fill=white,text width=5.85cm] (expansion) at (3.65,3.45)
    {\(\displaystyle x(h)=p+\frac{h^q}{\Gamma(q+1)}f(p)\)\\[-0.1em]
     \(\displaystyle\hphantom{x(h)={}}+
       \frac{h^{2q}}{\Gamma(2q+1)}J_f(p)f(p)+O(h^{3q})\)\\[0.25em]
     la condición anula el coeficiente de \(h^{2q}\)};
  \draw[arr] (embed.south) to[bend left=12] (expansion.east);
  \node[draw=Rojo,rounded corners=2pt,fill=white,text=Rojo,align=center] at (3.65,1.55)
    {propone una salida local; no define\\una aceleración global ni un invariante};

  % Carril historico FPP-H.
  \path[draw=Naranja,very thick,rounded corners=6pt,fill=NaranjaClaro]
    (7.45,0.95) rectangle (14.85,7.00);
  \node[font=\bfseries,Naranja] at (11.15,6.67)
    {FPP--H: evento sobre una historia};
  \draw[GrisOscuro,very thick] (8.05,5.45)--(14.10,5.45);
  \foreach \x in {8.05,8.80,9.55,10.30,11.05,11.80,12.55,13.30,14.10}
    {\fill[GrisOscuro] (\x,5.45) circle (1.15pt);}
  \node[below] at (8.05,5.37) {\(0\)};
  \node[below] at (14.10,5.37) {\(t\)};
  \node[above] at (11.05,5.55) {registro \(x|_{[0,t]}\)};

  \node[box,fill=white,text width=2.35cm] (profile) at (9.25,3.72)
    {perfil \(m_t\)\\dato para continuar};
  \node[box,fill=white,text width=3.10cm] (operator) at (12.92,3.72)
    {\(\mathcal A_{q,q}[x](t)\)\\usa el registro histórico};
  \draw[arr] (9.20,5.40)--(profile.north);
  \draw[arr] (12.65,5.40)--(operator.north);
  \node[text=Rojo,font=\large] at (11.05,3.72) {\(\not\equiv\)};
  \node[text=Rojo,fill=NaranjaClaro,inner sep=1pt] at (11.15,4.45)
    {objetos distintos};
  \node[box,fill=white,text width=3.10cm] (event) at (12.92,1.85)
    {evento: \(\mathcal A_{q,q}[x](t)=0\)\\con \(f(x(t))\neq0\)};
  \draw[arr] (operator)--(event);
  \node[box,fill=white,text width=2.35cm] at (9.25,1.85)
    {la trayectoria portadora y su memoria deben conservarse};

  \node[box,fill=Gris] at (7.50,0.30)
    {sólo para \(q=1\), y bajo regularidad, la regla de cadena recupera
     \(\frac{\mathrm d}{\mathrm dt}f(x(t))=J_f(x(t))f(x(t))\)};
\end{tikzpicture}%
}
\par\smallskip
\begin{minipage}{0.92\linewidth}
\small\textbf{Dos objetos diferentes para \(0<q<1\).}
FPP--A cancela el primer término de curvatura no trivial de la expansión local
Picard--Caputo y sirve como semilla. FPP--H es un evento dependiente del pasado
a lo largo de una trayectoria. Ninguno de los dos, por sí solo, certifica una
cuenca, recurrencia ni un atractor.
\end{minipage}
\end{center}


En campos afines globales, la linealidad de Caputo hace coincidir las
expresiones \(\mathcal A_{q,q}[x](t)\) y \(J_ff(x(t))\). En Chua PWL,
esa coincidencia es local mientras toda la memoria relevante permanezca
en una misma región; después de cruzar una frontera, FPP--H conserva
contribuciones de las regiones anteriores.

\begin{table}[ht]
\centering
\small
\caption{Interpretación geométrica de FPP--A y FPP--H.}
\label{tab:geofo-fpp-geometry}
\begin{tabularx}{\linewidth}{@{}L{0.16\linewidth}Y Y@{}}
\toprule
Objeto & Espacio natural & Qué puede afirmar \\
\midrule
FPP--A & \(p\in\mathbb R^n\), embebido como \(\iota(p)\in\mathfrak C\) &
Cancela la primera curvatura no trivial del arranque Picard--Caputo; propone
una semilla. No certifica un evento posterior ni una cuenca.\\
FPP--H & Trayectoria portadora más historia de convolución &
Detecta un mínimo o cero de aceleración secuencial dependiente de memoria. Su
proyección no es un punto perpetuo independiente del pasado.\\
PP clásico & \(p\in\mathbb R^n\) para \(q=1\) &
La regla de cadena reduce aceleración a \(J_ff\); FPP--A y FPP--H coinciden en
el límite entero bajo regularidad.\\
\bottomrule
\end{tabularx}
\end{table}

\fi

\ifnum\HAFOCaputoShowProtocol=1\relax
\ifdefined\HAFOCaputoBlockSelect
\section{Protocolo geométrico para sistemas fraccionarios}
\fi
\subsection{Protocolo numérico geométrico para Caputo}
\label{subsec:geofo-numerical-protocol}

El protocolo C1--C10 extiende los niveles L0--L6 mediante controles
funcionales de memoria.

\begin{enumerate}[label=\textbf{C\arabic*.}]
 \item \textbf{Fijar el operador.} Registrar definición de Caputo, órdenes
 \(q_j\), terminal inferior, datos iniciales, kernel, malla, integrador y
 política de memoria. Una ventana truncada o un reinicio define otra
 especificación numérica.

 \item \textbf{Probar existencia en el dominio usado.} Documentar Lipschitz
 global, crecimiento lineal, región absorbente o, como mínimo, el intervalo
 máximo y los controles de no explosión. Separar esta prueba de la mera
 finitud numérica.

 \item \textbf{Levantar equilibrios.} Resolver \(f(e)=0\), formar
 \(E_e=\iota(e)\), calcular \(J_f(e)\) y el margen de Matignon. Esto fija la
 geometría local desde la cual se diseñarán perturbaciones.

 \item \textbf{Comparar reinicio y continuación heredada.} Para cada estado candidato \(p\),
 ejecutar (a) un PVI reiniciado desde \(\iota(p)\) y (b), cuando \(p\) provenga
 de una trayectoria, la continuación desde su perfil heredado \(m_t\). No
 confundir una evolución con la otra.

 \item \textbf{Aproximar el perfil, no sólo el punto.} En tiempos seleccionados
 \(t_n\), reconstruir en una malla \(\theta_\ell\geq0\)
 \begin{equation}
  m_{t_n}(\theta_\ell)\approx x_0+
  \frac1{\Gamma(q)}\sum_{j=0}^{n-1}
  \int_{t_j}^{t_{j+1}}
  (t_n+\theta_\ell-s)^{q-1}\widehat f_j(s)\,\mathrm ds,
  \label{eq:geofo-discrete-profile}
 \end{equation}
 donde \(\widehat f_j\) es la misma interpolación usada por el integrador.
 Comparar estos perfiles al refinar \(h\), no combinar cuadraturas
 incompatibles.

 \item \textbf{Diagnosticar recurrencia en dos métricas.} Informar
 \(R_x\) y una aproximación truncada de \(R_{\mathfrak C}\) sobre varias
 ventanas \(\theta\in[0,\Theta]\). La convergencia sólo en \(\mathbb R^n\) se
 etiqueta como recurrencia proyectada.

 \item \textbf{Construir retornos funcionales.} Detectar cruces orientados por
 interpolación y cambio de signo. Conservar, junto a cada punto de sección, un
 descriptor de memoria. Probar si el retorno proyectado es univaluado:
 perfiles cercanos en estado pero separados en memoria no deben producir
 salidas incompatibles si se quiere ajustar un mapa finito-dimensional.

 \item \textbf{Sondear la cuenca de reinicio.} Muestrear bolas o superficies
 alrededor de cada equilibrio, iniciar siempre con el mismo perfil constante y
 usar el clasificador preregistrado. Reportar el primer radio de contacto y los
 casos no clasificados.

 \item \textbf{Sondear direcciones funcionales.} Elegir una familia finita y
 declarada
 \begin{equation}
  m_{e,\alpha}(\theta)=E_e(\theta)+
  \sum_{k=1}^{K}\alpha_k\psi_k(\theta),
  \qquad \|\alpha\|\leq\delta,
  \label{eq:geofo-history-probes}
 \end{equation}
 con funciones base continuas \(\psi_k\). Si se pretende estudiar ocultedad
 funcional, demostrar primero que los perfiles pertenecen al espacio
 \(X_{\mathrm{adm}}\) declarado e integrar su continuación sin abandonar
 ese espacio. Sin esa justificación, son sondeos de la extensión de Volterra,
 útiles para estudiar sensibilidad al término libre, pero no pruebas sobre
 memorias admisibles. La atracción abierta en todo \(\mathfrak C\) ya está
 descartada analíticamente por \cref{prop:geofo-full-space-obstruction}.

 \item \textbf{Separar niveles de conclusión.} Exigir, en orden: identidad
 algebraica; convergencia del integrador; conjunto acotado persistente;
 recurrencia funcional; atracción de una familia local \(\mathcal D_A\)
 de reinicios o de perfiles admisibles; separación de cuenca de reinicio; y,
 sólo en un \(X_{\mathrm{adm}}\) construido y con argumentos adicionales,
 ocultedad funcional o caos topológico.
\end{enumerate}

Para órdenes inconmensurables \(q_j\), el perfil se define componente a
componente,
\begin{equation}
 [m_t(\theta)]_j=x_{0,j}+\frac1{\Gamma(q_j)}
 \int_0^t(t+\theta-s)^{q_j-1}f_j(x(s))\,\mathrm ds.
 \label{eq:geofo-incommensurate-profile}
\end{equation}
Cada componente requiere su propio orden y su correspondiente conjunto de
pesos.

\subsection{Ejemplo: relajación escalar}

Considérese
\begin{equation}
 {}^{\mathrm C}D_{0+}^{q}x=-a x,
 \qquad a>0,\quad 0<q<1.
 \label{eq:geofo-relaxation}
\end{equation}
El único equilibrio es \(e=0\) y
\begin{equation}
 x(t)=x_0E_q(-at^q)\longrightarrow0.
 \label{eq:geofo-relaxation-solution}
\end{equation}
El criterio de Matignon se cumple porque
\(|\arg(-a)|=\pi>q\pi/2\). El perfil en el tiempo \(t\) es
\begin{equation}
 m_t(\theta)=x_0-
 \frac a{\Gamma(q)}\int_0^t
 (t+\theta-s)^{q-1}x_0E_q(-as^q)\,\mathrm ds.
 \label{eq:geofo-relaxation-profile}
\end{equation}
Al poner \(\theta=0\) se recupera \(m_t(0)=x(t)\), pero para
\(\theta>0\) el perfil no es constante. Por eso continuar y reiniciar desde
\(x(t)\) generan curvas distintas, aunque ambas converjan a cero.

Este ejemplo permite identificar todos los objetos:
\begin{itemize}
 \item el único perfil constante de equilibrio es \(E_0(\theta)\equiv0\);
 \item la órbita verdadera es \(\{m_t:t\geq0\}\subset\mathfrak C\);
 \item la curva \(x(t)\) es su proyección por \(\operatorname{ev}_0\);
 \item toda la recta pertenece a la cuenca de reinicio de \(E_0\);
 \item las aplicaciones puntuales no forman un semigrupo por
 \cref{eq:geofo-false-ml-product};
 \item ningún reinicio produce un ciclo periódico no constante ni un
 destino oculto diferente de \(E_0\).
\end{itemize}
La \cref{prop:geofo-reset-equilibrium-separation} asegura la convergencia
de cada perfil a \(E_0\). En este caso, la estimación de su demostración es
\[
 \sup_{0\leq\theta\leq\Theta}|m_t(\theta)|
 \leq\left(1+\frac{a\Theta^q}{\Gamma(q+1)}\right)
       \sup_{u\in[t,t+\Theta]}|x(u)|\longrightarrow0.
\]
La cota es uniforme para \(x_0\) en conjuntos acotados, de modo que
\(\{E_0\}\) atrae el universo \(\mathcal D\). Sin embargo, no atrae
ninguna vecindad abierta de \(E_0\) en todo \(\mathfrak C\): cada una
contiene puntos fijos artificiales
\cref{eq:geofo-artificial-fixed-profile} con \(c\ne0\).
La relajación escalar constituye un caso de verificación para comprobar que
una implementación distingue continuación y reinicio antes de estudiar Chua.

\subsection{Alcance de los resultados geométricos}

\begin{longtable}{@{}L{0.22\linewidth}L{0.34\linewidth}L{0.36\linewidth}@{}}
\caption{Resultados y condiciones del análisis geométrico de Caputo.}
\label{tab:geofo-theory-boundary}\\
\toprule
Tipo de resultado & Contenido & Fundamento o condición \\
\midrule
\endfirsthead
\toprule
Tipo de resultado & Contenido & Fundamento o condición \\
\midrule
\endhead
Teoría establecida & Equivalencia Caputo--Volterra; ausencia general de
semigrupo sobre \(\mathbb R^n\); semidinámica sobre \(\mathfrak C\); criterio
de Matignon para reinicios; no periodicidad exacta de soluciones no constantes
del PVI bajo hipótesis finitas. &
Respaldado por
\cite{Matignon1996,CongTuan2017Nonlocal,DoanKloeden2021Semidynamical,
TavazoeiHaeri2009,KaslikSivasundaram2012}.\\
Consecuencia directa & Perfil \(m_t\), separación algebraica de memoria,
identidades de continuar frente a reiniciar, proyección
\(\operatorname{ev}_0\), operador de retorno funcional y obstrucción de
atracción abierta en todo \(\mathfrak C\). &
Derivación algebraica de los perfiles y
\cref{prop:geofo-full-space-obstruction}.\\
Resultado para Lorenz & Existencia global y absorción de reinicios;
precompacidad de sus perfiles para cada familia acotada. &
Demostrado para los parámetros indicados en
\cref{prop:geofo-lorenz-boundedness,prop:geofo-profile-precompact}; no
certifica un destino caótico u oculto.\\
Atracción y cuencas & Jerarquía entre cuenca funcional y cuenca de
reinicio con hipótesis de admisibilidad; atracción relativa a
\(\mathcal D_A\); ocultedad en \(X_{\mathrm{adm}}\); protocolo C1--C10;
interpretación geométrica complementaria de función descriptiva, continuación y
FPP--A/FPP--H. &
Formulación relativa que requiere construir el espacio admisible y demostrar
la atracción de las familias declaradas.\\
Problemas de los modelos & Identificación y atracción de candidatos concretos;
regiones absorbentes para RF, perfiles admisibles cerca de
equilibrios y convergencia de retornos con memoria. &
No se deduce de la evidencia finita presentada.\\
Estructuras globales & Variedades estables/inestables intrínsecas,
homoclinicidad, índice de Conley, descomposición de Morse, entropía topológica,
herraduras y geometría fractal de cuencas en el espacio funcional. &
Requieren un espacio de fases apropiado y argumentos globales de invariancia,
aislamiento o intersección transversal.\\
\bottomrule
\end{longtable}

\subsection{Conexión con la función descriptiva y la continuación causal}

La función descriptiva entera o de Machado genera una amplitud, frecuencia y
semilla puntual. En términos geométricos, esa semilla entra por
\(\iota(x_{\mathrm{sem}})\), es decir, por la sección de reinicio. El predictor
armónico basado en \((\mathrm i\omega)^q\) describe una respuesta estacionaria
tipo Liouville--Weyl; no produce por sí mismo una órbita periódica exacta del
PVI Caputo con terminal inferior finito.

La continuación causal de \cref{sec:continuacion_causal} sí conserva una sola
historia. Durante la rampa \(\varepsilon(t)\), el campo es no autónomo, de modo
que el proceso de localización no es una órbita invariante del sistema objetivo.
Sólo después de fijar \(\varepsilon=1\) puede analizarse la cola respecto de la
semidinámica autónoma final. Reiniciar en cada valor de \(\varepsilon\) define
una familia de perfiles constantes distinta y no es topológicamente equivalente
a recorrer una sola curva en \(\mathfrak C\).

Los métodos se complementan así:
\begin{equation}
 \begin{array}{c}
 \text{función descriptiva o FPP--A}\\[-0.1em]
 \text{propone }x_{\mathrm{sem}}\in\mathbb R^n
 \end{array}
 \xrightarrow{\ \iota\ }
 \mathfrak C
 \xrightarrow{\ \text{continuación causal}\ }
 m_T
 \xrightarrow{\ S_t\ }
 \text{dinámica objetivo},
 \label{eq:geofo-method-chain}
\end{equation}
mientras FPP--H actúa como diagnóstico a lo largo de la historia y el sondeo de
cuenca clasifica el destino bajo una convención de memoria explícita.

\subsection{Aplicación a los sistemas fraccionarios analizados}

\subsubsection{Chua PWL y Chua arctangente}

La saturación continua PWL es globalmente Lipschitz y el campo de Chua tiene
crecimiento a lo sumo lineal; la arctangente también es globalmente Lipschitz.
Estos casos son, por ello, los más adecuados para construir una semidinámica
global sobre \(\mathfrak C\).

Los planos de conmutación \(x=\pm1\) siguen siendo hipersuperficies en la
proyección puntual, pero no dividen el espacio funcional en sólo tres regiones:
un perfil cuyo valor actual está en la región central puede conservar memoria
de visitas a las regiones externas. Esto explica por qué una identidad PWA
local no puede aplicarse retrospectivamente a toda la convolución.

Los equilibrios \(e_j\) se levantan a perfiles fijos \(E_{e_j}\). Los sondeos
esféricos realizados alrededor de \(e_j\) exploran únicamente
\(\iota[B(e_j,\rho)]\). Para fortalecer la afirmación topológica habría que
construir primero un espacio \(X_{\mathrm{adm}}\) adecuado y después
estudiar perfiles no constantes próximos a \(E_{e_j}\) dentro de él.
La suavidad y la existencia global sobre \(\mathfrak C\) no eliminan
la obstrucción de \cref{prop:geofo-full-space-obstruction}.

En Chua arctangente, la suavidad permite calcular sin ambigüedad \(J_f\), la
linealización de Matignon y FPP--A. Aun así, la semilla sesgada y el cierre de
Machado siguen siendo predictores en la sección de reinicio; sus residuos no
demuestran invariancia funcional.

\subsubsection{Lorenz generalizado y Rabinovich--Fabrikant}

Los campos de Lorenz y RF son suaves, pero polinomiales y no globalmente
Lipschitz. En el Lorenz generalizado de \cref{part:nonchua}, la estructura
cuadrática permite cerrar la existencia global y la absorción para reinicios.

\begin{proposition}[existencia global y absorción del Lorenz generalizado]
\label{prop:geofo-lorenz-boundedness}
Sean \(\sigma=3.4\), \(r=6.8\), \(a=-0.5\), un orden conmensurable
\(0<q\leq1\) y un terminal inferior fijo, tomado como cero. Para el campo
\begin{equation}
 F(x,y,z)=\begin{pmatrix}
 -3.4(x-y)+0.5yz\\
 6.8x-y-xz\\
 -z+xy
 \end{pmatrix},
 \qquad
 V(x,y,z)=x^2+y^2+\tfrac12(z-20.4)^2,
 \label{eq:geofo-lorenz-lyapunov}
\end{equation}
todo PVI de Caputo con reinicio \(X(0)=X_0\in\mathbb R^3\) tiene una
solución global única y satisface
\begin{equation}
 V(X(t))\leq V(X_0)E_q(-t^q)
       +208.08\,[1-E_q(-t^q)].
 \label{eq:geofo-lorenz-comparison}
\end{equation}
En particular,
\begin{equation}
 V(X(t))\leq\max\{V(X_0),208.08\},
 \qquad \limsup_{t\to\infty}V(X(t))\leq208.08.
 \label{eq:geofo-lorenz-ultimate-bound}
\end{equation}
Para cada \(\varepsilon>0\), el elipsoide
\(\{X:V(X)\leq208.08+\varepsilon\}\) absorbe las trayectorias de
cualquier familia acotada de reinicios, con tiempo de entrada uniforme
para esa familia.
\end{proposition}

\begin{proof}
Primero supóngase \(0<q<1\). La suavidad del campo da existencia y unicidad
local. En cualquier intervalo
compacto anterior al tiempo maximal, la solución pertenece a la clase
\(X-X_0=I^qF(X)\), con \(F(X)\) acotado. La desigualdad de composición
para funciones convexas en esta clase de soluciones
\cite[Lema~3.1]{Gomoyunov2018Convex} se aplica a
\(\xi=X-X_0\) y \(W(\xi)=V(X_0+\xi)-V(X_0)\): así,
\(\xi(0)=0\), \(W(0)=0\), y la Hessiana constante
\(D^2W=\operatorname{diag}(2,2,1)\) es positiva definida.
Por tanto, \(W\) es convexa y \(\nabla W\) es Lipschitz.
El lema garantiza asimismo
\(V(X(\cdot))-V(X_0)\in I^q(L^\infty)\) en cada intervalo compacto;
la derivada y la posterior comparación escalar están definidas en esa clase.
Se obtiene
\({}^{\mathrm C}D^qV(X(t))\leq\nabla V(X(t))\cdot F(X(t))\)
para casi todo \(t\) del intervalo considerado.
No se utiliza una regla de cadena fraccionaria con igualdad. Para precisar
el signo en el caso suave, fije $t$ y ponga
\[
 \Phi(s)=V(X(s))-V(X(t))-\nabla V(X(t))\cdot(X(s)-X(t))\ge0.
\]
La convexidad da la desigualdad y $\Phi(t)=0$. La diferencia entre
${}^{\mathrm C}D^q[V(X)](t)$ y
$\nabla V(X(t))\cdot{}^{\mathrm C}D^qX(t)$, multiplicada por
$\Gamma(1-q)$, es
\[
 \int_0^t(t-s)^{-q}\Phi'(s)\,\mathrm ds
 =-t^{-q}\Phi(0)-q\int_0^t(t-s)^{-q-1}\Phi(s)\,\mathrm ds\le0.
\]
La integración por partes se realiza primero hasta $t-\varepsilon$.
Si $X$ y $\nabla V$ son Lipschitz, $\Phi(s)=O((t-s)^2)$;
el término de extremo tiende a cero y la integral converge.
El lema citado extiende el argumento por aproximación a la clase
$I^q(L^\infty)$ usada aquí, sin exigir esa regularidad ordinaria de la
trayectoria. El cálculo algebraico es
\begin{align}
 \nabla V\cdot F
 &=2x[-3.4(x-y)+0.5yz]
   +2y[6.8x-y-xz]\nonumber\\
 &\quad +(z-20.4)(-z+xy)\nonumber\\
 &=-6.8x^2-2y^2-z^2+20.4z\nonumber\\
 &=-V+208.08-5.8x^2-y^2-\tfrac12z^2
 \leq -V+208.08.
 \label{eq:geofo-lorenz-dissipation}
\end{align}
Los términos cúbicos suman \((1-2+1)xyz=0\), y el coeficiente de
\(xy\) es \(6.8+13.6-20.4=0\). Además,
\(\tfrac12(20.4)^2=208.08\), lo que da la última igualdad.

Para detallar la comparación escalar, sea \(v(t)=V(X(t))\) y escriba
\(r(t)=208.08-v(t)-{}^{\mathrm C}D^qv(t)\geq0\) casi en todas partes.
El residuo \(r\) pertenece a \(L^\infty\) en cada horizonte finito.
La fórmula de variación de constantes de
\({}^{\mathrm C}D^qv=-v+208.08-r\), obtenida al resolver su ecuación
integral lineal, es
\begin{align*}
 v(t)&=v(0)E_q(-t^q)+208.08[1-E_q(-t^q)]\\
 &\quad-\int_0^t(t-s)^{q-1}E_{q,q}(-(t-s)^q)r(s)\,\mathrm ds.
\end{align*}
El término constante se integra usando
\(\int_0^t u^{q-1}E_{q,q}(-u^q)\,\mathrm du=1-E_q(-t^q)\).
Esta identidad y
\[
 t^{q-1}E_{q,q}(-t^q)
 =-\frac{\mathrm d}{\mathrm dt}E_q(-t^q)
\]
se deducen al integrar y derivar la serie de Mittag--Leffler, respectivamente.
Para \(0<q<1\), \(E_q(-t^q)\) es positiva y decreciente
\cite{Podlubny1999}; el núcleo de la integral sustraída es, por tanto,
no negativo. Se obtiene
\cref{eq:geofo-lorenz-comparison}. Puesto que
\(0<E_q(-t^q)\leq1\) y \(E_q(-t^q)\to0\), se obtiene
\cref{eq:geofo-lorenz-ultimate-bound} mientras exista la solución.

La función \(V\) es coerciva: sus subniveles son compactos. La cota anterior
mantiene la solución maximal en uno de ellos, \(K\). Para excluir un tiempo
maximal finito \(T_*\), elija un corte suave igual a uno en una vecindad de
\(K\) y de soporte compacto. El campo recortado es globalmente Lipschitz y
su PVI de Volterra tiene solución global. Por unicidad, ésta coincide con la
original en \([0,T_*)\). Por continuidad, su estado en \(T_*\) pertenece
a \(K\), y permanece en la vecindad donde ambos campos coinciden durante
un intervalo posterior. La misma solución integral, con toda su historia,
prolonga entonces el PVI original más allá de \(T_*\), contradicción.
La solución es global. Si \(B\) es un
conjunto acotado de reinicios, \(C_B=\sup_{X_0\in B}V(X_0)<\infty\), y
\[
 \sup_{X_0\in B}V(X(t;X_0))
 \leq208.08+(C_B-208.08)_+E_q(-t^q).
\]
El último término es menor que \(\varepsilon\) para todo tiempo
suficientemente grande, uniformemente en \(B\). Para \(q=1\) se usa
la desigualdad ordinaria y \(E_1(-t)=e^{-t}\).
\end{proof}

\begin{proposition}[precompacidad de perfiles de reinicios acotados]
\label{prop:geofo-profile-precompact}
Sea una familia de soluciones globales con reinicio de orden
\(0<q<1\), cuyos estados y campos satisfagan uniformemente
\(\|X(t)\|\leq R\) y \(\|F(X(t))\|\leq M\) para \(t\geq0\).
Entonces sus perfiles heredados verifican, para
\(0\leq\theta_1\leq\theta_2\),
\begin{align}
 \|m_t(\theta_2)-m_t(\theta_1)\|
 &\leq\frac{M}{\Gamma(q+1)}|\theta_2-\theta_1|^q,
 \label{eq:geofo-profile-holder}\\
 \|m_t(\theta)\|&\leq R+\frac{M\theta^q}{\Gamma(q+1)}.
 \label{eq:geofo-profile-bound}
\end{align}
La unión de estas órbitas de perfiles es precompacta en
\(\mathfrak C\). En particular, el resultado se aplica al Lorenz anterior
para cada conjunto acotado de reinicios, incluido \(q=0.995\).
\end{proposition}

\begin{proof}
La monotonía de \(s^{q-1}\) permite acotar la diferencia de las integrales
que definen \(m_t\) por
\[
 \frac{M}{\Gamma(q+1)}
 \bigl[(t+\theta_1)^q-\theta_1^q
       -(t+\theta_2)^q+\theta_2^q\bigr].
\]
Como \((t+\theta_1)^q\le(t+\theta_2)^q\), el corchete es a lo sumo
\(\theta_2^q-\theta_1^q\). La subaditividad de \(s^q\) da
\(\theta_2^q=(\theta_1+\theta_2-\theta_1)^q
\le\theta_1^q+(\theta_2-\theta_1)^q\), de donde resulta
\cref{eq:geofo-profile-holder}.
Como \(m_t(0)=X(t)\), se obtiene también
\cref{eq:geofo-profile-bound}. En cada intervalo \([0,\Theta]\), los
perfiles están uniformemente acotados y son equicontinuos. El teorema de
Arzelà--Ascoli proporciona una subsucesión convergente en \([0,1]\), de
ella otra en \([0,2]\), y así sucesivamente. La diagonal converge en cada
compacto a un perfil continuo; los límites coinciden en los intervalos
superpuestos. Para la métrica \cref{eq:geofo-compact-open-metric}, se fija
primero \(N\) con \(2^{-N}<\varepsilon/2\) y después se hace menor que
\(\varepsilon/2\) la suma de los primeros \(N\) términos mediante esa
convergencia uniforme. La cola es menor que \(2^{-N}\). Así, toda sucesión
tiene una subsucesión convergente en \(\mathfrak C\), lo que demuestra
precompacidad.
Para Lorenz, \cref{prop:geofo-lorenz-boundedness} proporciona un compacto
común para los estados de la familia, y la continuidad de \(F\) da \(M\).
\end{proof}

Para interpretar el cierre como espacio dinámico no se invoca la proposición
global de \(\mathfrak C\) directamente con el campo polinomial. Fijado
un conjunto acotado de reinicios \(B\), se toma una extensión globalmente
Lipschitz \(\widetilde F\) que coincide con \(F\) en un entorno del
compacto común visitado; un corte suave de soporte compacto la proporciona.
El semiflujo de \(\widetilde F\) coincide con el original en esas órbitas.
Su continuidad muestra que el cierre de la unión de órbitas es compacto y
positivamente invariante. En cualquier perfil de ese cierre, los estados de
su continuación permanecen en el mismo compacto, por continuidad de
\(\operatorname{ev}_0S_t\); por tanto, allí también coinciden ambos campos.
Se obtiene así un semiflujo continuo del Lorenz original sobre ese cierre.
Cada órbita tiene un conjunto \(\omega\)-límite no vacío y compacto.
Su invariancia positiva se sigue de la continuidad. Para la inclusión
inversa, si \(a=\lim_jS_{t_j}m\) y \(s>0\), la compacidad permite
extraer una subsucesión \(S_{t_j-s}m\to b\in\omega(m)\); entonces
\(S_sb=a\). Así, \(S_s\omega(m)=\omega(m)\) sin suponer inversas
del semiflujo. También atrae la órbita: una sucesión de tiempos con distancia
positiva uniforme a \(\omega(m)\) tendría, por compacidad, un límite
perteneciente al propio conjunto, contradicción.
Esta construcción depende de \(B\):
no demuestra atracción de un abierto de todo \(\mathfrak C\), no
identifica un mismo destino para diferentes reinicios y no establece
ocultedad funcional ni caos.

\ifHAFOPedagogy
\begin{bloqueestudio}{Ejercicios}
Compruebe las cancelaciones de \cref{eq:geofo-lorenz-dissipation} y deduzca
cotas explícitas para \(|x|\), \(|y|\) y \(|z-20.4|\) a partir de
\(V\leq C\). Para un conjunto de reinicios con \(V(X_0)\leq C_B\),
exprese la condición que debe satisfacer el tiempo de absorción en términos
de \(E_q(-t^q)\) y \(\varepsilon\). Justifique cada uso de
Arzelà--Ascoli en \cref{prop:geofo-profile-precompact} y explique por qué
estas cotas no prueban la existencia de un atractor oculto concreto.
\end{bloqueestudio}
\fi

En Rabinovich--Fabrikant permanecen pendientes una cota que excluya explosión
en el dominio pertinente y el control correspondiente de perfiles. La
estimación de Lorenz no se transfiere a RF sin una demostración propia.
Los cálculos de equilibrios, espectros y realizaciones de Lur'e de
\cref{part:nonchua}, así como las semillas FPP--A, deben complementarse
con reproducción Caputo y estudio de cuenca para identificar los destinos
dinámicos de ambos modelos.

La simetría puntual de Lorenz o RF puede levantar una simetría funcional
mediante aplicación punto a punto si conmuta con el campo y con el kernel. Esa
simetría relaciona perfiles y destinos, pero no demuestra por sí sola que dos
nubes simétricas sean atractores distintos; puede tratarse de proyecciones de
un mismo conjunto funcional simétrico.

\subsubsection{Muñoz--Pacheco}

El sistema de Muñoz--Pacheco es valioso porque no posee equilibrios para los
parámetros considerados \cite{munoz2018}. Sus cuatro raíces FPP--A verificadas
proporcionan estados de arranque algebraicamente motivados. El cálculo
presentado demuestra integración finita desde esas semillas, pero no certifica:
\begin{itemize}
 \item un compacto invariante con atracción de una familia local
 \(\mathcal D_A\) de reinicios;
 \item convergencia de las cuatro memorias hacia el mismo conjunto;
 \item un evento FPP--H estable bajo refinamiento;
 \item una propiedad topológica de caos.
\end{itemize}
En ausencia de equilibrios, el análisis geométrico pertinente consiste en
construir y comparar cuencas de diferentes perfiles de reinicio, en lugar de
aplicar de forma vacía el criterio clásico de vecindades de equilibrios.

\begin{table}[ht]
\centering
\small
\caption{Objeto geométrico y verificación requerida por familia.}
\label{tab:geofo-case-program}
\begin{tabularx}{\linewidth}{@{}L{0.20\linewidth}L{0.31\linewidth}Y@{}}
\toprule
Familia & Resultados disponibles & Problema geométrico \\
\midrule
Chua PWL fraccionario & Perfiles de reinicio, continuación causal, sondeos
puntuales y regiones PWA & Espacio de perfiles admisibles, atracción relativa
y retornos con memoria.\\
Chua arctangente & Campo suave, Matignon, semillas sesgadas, Machado y FPP--A &
Comparar continuación heredada con reinicio y estudiar cuencas en el espacio
admisible declarado.\\
Lorenz \(q=0.995\) & Álgebra, equilibrios, espectros, FPP--A; existencia
global, absorción y precompacidad de perfiles de reinicios acotados &
Reproducción dinámica, identificación del destino y atracción de
\(\mathcal D_A\), cuenca y ocultedad de reinicio.\\
RF \(q=0.998\) & Álgebra, cinco equilibrios, espectros y ocho FPP--A &
Existencia global y control de perfiles; reproducción y cuenca; deduplicación
por simetría en el espacio funcional.\\
Muñoz--Pacheco \(q=0.97\) & Cuatro FPP--A e integración PECE finita con historia completa &
Convergencia temporal/espacial, destino común y búsqueda FPP--H refinada.\\
\bottomrule
\end{tabularx}
\end{table}

\subsection{Reformulación del problema geométrico en sistemas de Caputo}

En orden entero se pregunta dónde está la órbita en \(\mathbb R^n\), qué
variedades la organizan y qué cuenca la conduce al atractor. En Caputo deben
formularse dos preguntas simultáneas:
\begin{equation}
 \text{¿qué se observa en }\mathbb R^n?
 \qquad\text{y}\qquad
 \text{¿qué perfil evoluciona en }\mathfrak C?
 \label{eq:geofo-two-geometric-questions}
\end{equation}
La primera produce los retratos de fase y es indispensable para interpretar el
circuito o modelo físico. La segunda exige declarar el espacio y la familia
de datos antes de estudiar invariancia, retorno, atracción y ocultedad en
sentido dinámico. La evaluación
\(\operatorname{ev}_0\) relaciona el estado funcional con su observación
puntual, mientras \(\iota\) incorpora las condiciones iniciales convencionales.

Una formulación funcional conserva las semillas de función descriptiva y
FPP--A; la continuación debe preservar
memoria; FPP--H se registra como evento de historia; y las cuencas puntuales se
etiquetan como secciones de reinicio. Sólo después de estas traducciones tiene
sentido buscar variedades, índices topológicos, entropía o mecanismos
homoclínicos para los sistemas fraccionarios concretos.
\fi
